\documentclass[11pt,onecolumn]{article}
\usepackage[margin=1in]{geometry}
\usepackage{amsmath,amssymb,amsthm}
\usepackage{booktabs}
\usepackage{graphicx}
\usepackage{xcolor}
\usepackage[numbers,sort&compress]{natbib}
\usepackage{hyperref}
\usepackage{enumitem}
\usepackage{caption}
\usepackage{subcaption}
\usepackage{array}
\usepackage{colortbl}
\usepackage{multirow}
\usepackage{float}
\usepackage{parskip}

\usepackage{setspace}

\definecolor{deepblue}{RGB}{46,64,87}
\definecolor{rowhi}{RGB}{255,242,204}
\definecolor{rownull}{RGB}{234,244,234}
\definecolor{cat_env}{RGB}{33,150,243}
\definecolor{cat_eng}{RGB}{255,87,34}
\definecolor{cat_spa}{RGB}{76,175,80}
\definecolor{cat_tmp}{RGB}{156,39,176}
\definecolor{cat_sta}{RGB}{96,125,139}

\hypersetup{colorlinks=true, linkcolor=deepblue, citecolor=deepblue, urlcolor=deepblue}

\newtheorem{proposition}{Proposition}
\newtheorem{corollary}{Corollary}[proposition]
\newtheorem{lemma}{Lemma}
\newtheorem{remark}{Remark}
\newtheorem{definition}{Definition}

\title{%
  \textbf{\color{deepblue}Non-Commutativity of Sobol Sensitivity Indices\\
  Under Output Transformations of Functional, Spatial, and Scalar Models}\\[0.4em]
  \large A Systematic Empirical Sweep Across 33 Transforms, 10 Benchmarks,
  and Four Transform Categories
}
\author{Dylan Hettinger \\ \small Research Draft --- March 2026 \\ \small \textit{Prepared for submission to the Electronic Journal of Statistics}}
\date{}

\begin{document}
\maketitle

\begin{abstract}
Variance-based (Sobol) global sensitivity analysis is routinely applied to model
outputs that have been transformed, aggregated, thresholded, or otherwise
post-processed for physical or statistical reasons.
A widespread implicit assumption is that sensitivity indices of the transformed
output can be inferred from those of the raw output.
We demonstrate both analytically and empirically---across a registry of 33
physically motivated transforms applied to 10 benchmark functions---that this
assumption fails for virtually every nonlinear transformation, with direct
consequences for factor-fixing and model-reduction decisions.

We introduce \emph{functional sensitivity analysis under output operators}
(FSAO) as a unifying framework: Sobol indices are defined pointwise over a
discretised functional output (spatial grid, temporal trajectory, or scalar),
and the question of non-commutativity is whether $S_i(T \circ Y) = S_i(Y)$
for a post-processing operator $T$.  We prove that $T$ commutes with Sobol
index computation if and only if $T$ is affine (Proposition~1), and that no
nonlinear $T$ commutes in general.  Two further propositions, motivated by
empirical findings and proved within the Results section, characterise the two
practically important special cases: the change-of-support (COS) operator
commutes if and only if the underlying sensitivity structure is spatially
homogeneous (Proposition~2), and all strictly monotone transforms produce
identical Decision Score $D$ and Sensitivity Shift $\Delta$ on scalar benchmarks
(Proposition~3).  All three results extend formally to total-effect Sobol indices
(Appendix~A).

The empirical sweep covers 5 transform categories (environmental, engineering,
spatial, temporal, statistical) applied to 2 spatial benchmarks
(Campbell2D, Campbell3D), 2 temporal/functional benchmarks
(Damped Oscillator, Lotka-Volterra), and 6 scalar benchmarks
(Ishigami, SobolG, Borehole, Piston, WingWeight, OTL Circuit).
We score each of the 330 (benchmark, transform) pairs using a Decision Score
$D \in [0,1]$ (soft threshold margin around $\tau=0.05$) and a Sensitivity
Shift $\Delta \in [0,1]$ (Bray-Curtis dissimilarity), both purpose-designed
for bounded, cross-comparable measurement.
Environmental and engineering transforms produce $D = 0.36$--$0.37$ and
$\Delta = 0.60$--$0.61$ on average.
Temporal transforms, which are novel in this context, produce
$D = 0.18$ and $\Delta = 0.39$ on average, with linear temporal operators
(cumulative sum on the Damped Oscillator: $D = 0.064$) confirming the
commutativity theorem.
The Campbell3D benchmark, which we introduce here for the first time in the
sensitivity analysis literature, produces a $13.2\times$ amplification of
COS noncommutativity relative to Campbell2D at $8^3$ block scale.
All code and results are archived via \textit{sabench}~\citep{Hettinger2025}.
\end{abstract}

\tableofcontents
\clearpage

%% ─────────────────────────────────────────────────────────────────────────────
\section{Introduction}
\label{sec:intro}
%% ─────────────────────────────────────────────────────────────────────────────

Variance-based global sensitivity analysis (GSA) has become a standard
framework for model interrogation in the environmental, engineering, and earth
sciences \citep{Saltelli2008,IoossLemaitre2015,Razavi2021}. For a model
$Y = f(\mathbf{X})$ with independent uniform inputs, the Sobol first-order
index $S_i$ and total-effect index $S_i^T$ decompose output variance according
to individual and collective input contributions.

In practice, the quantity subjected to GSA is often not the raw model output.
Practitioners routinely apply transformations $T$ before or after sensitivity
analysis: log-transforming streamflow for flood-frequency analysis
\citep{Vogel1996}, applying depth-damage power laws for loss estimation
\citep{Merz2010}, computing exceedance probabilities for infrastructure design,
converting temperature fields to Arrhenius reaction rates for atmospheric
chemistry, or block-averaging distributed fields to match observational grids---the
classical change-of-support problem in spatial statistics
\citep{Cressie1993,GotwayYoung2002}. Temporal models are similarly post-processed: peak
response is extracted from oscillator outputs for structural design; cumulative
volume is integrated from flow-rate trajectories for pharmacokinetics; running
maxima define design loads in seismology.

The implicit assumption---universally unstated in applied papers that report
Sobol indices of transformed outputs---is that
$S_i(T \circ f) \approx S_i(f)$, or at least that relative rankings are preserved.
This assumption is false whenever $T$ is nonlinear.
Despite this, no systematic empirical study has quantified the magnitude and
practical significance of this violation across a broad, physically grounded
catalogue of transformations applied consistently to diverse benchmark types.
This paper closes that gap, introducing a unifying framework and comprehensive
empirical evidence.

\paragraph{Contributions.}
\begin{enumerate}
  \item A unified framework---\emph{functional sensitivity analysis under output
    operators} (FSAO)---that treats scalar, spatial, and temporal GSA as
    special cases of a common formalism, and operationalises the notion of
    commutativity between Sobol estimation and output post-processing
    (Section~\ref{sec:framework}).
  \item A complete, self-contained proof that linear (affine) operators are the
    unique class of transformations that commute with Sobol index computation
    (Proposition~\ref{prop:main} and Appendix~\ref{app:proofs}).
  \item A complete proof of the change-of-support (COS) theorem for Sobol
    indices, stated within the spatial results where the empirical evidence
    motivates it (Proposition~\ref{prop:cos}).
  \item A complete proof of scalar invariance under strictly monotone transforms,
    stated within the scalar results (Proposition~\ref{prop:scalar}).
  \item Formal extension of all three propositions to total-effect Sobol indices
    $S_i^T$, with proofs and numerical verification (Appendix~\ref{app:total_effects},
    Propositions~\ref{prop:main_T}--\ref{prop:scalar_T}), protecting the results
    from the objection that they are artefacts of the first-order decomposition.
  \item A fourth proposition characterising the rank/PIT transform as a
    near-safe variance-stabiliser that is ratio-preserving but not commutativity-preserving
    (Proposition~\ref{prop:rank}, Appendix~\ref{app:proofs}).
  \item A literature review connecting scalar Sobol analysis, pointwise spatial
    Sobol, functional output Sobol, and transformation commutativity
    (Section~\ref{sec:litreview}).
  \item A registry of 33 physically motivated transforms in 5 categories
    (environmental, engineering, spatial, temporal, statistical), together with
    purpose-designed bounded scoring metrics (Section~\ref{sec:setup}).
  \item The first sensitivity analysis of the Campbell3D benchmark
    (Section~\ref{sec:spatial_results}), and the first temporal sensitivity
    analysis under output transformations using established functional benchmarks
    (Section~\ref{sec:temporal_results}).
  \item Empirical confirmation that the cumulative-sum transform (linear)
    scores near zero ($D = 0.064$) on the Damped Oscillator, validating the
    commutativity theorem, while nonlinear temporal transforms
    (log-cumsum, exceedance duration, peak) score $D = 0.12$--$0.28$
    (Section~\ref{sec:temporal_results}).
\end{enumerate}

%% ─────────────────────────────────────────────────────────────────────────────
\section{Literature Review}
\label{sec:litreview}
%% ─────────────────────────────────────────────────────────────────────────────

\subsection{Scalar Sobol indices}
\label{sec:lit_scalar}

The theoretical foundations of variance-based sensitivity analysis rest on the
$L^2$ orthogonal decomposition of functions of independent random variables,
established by \citet{Hoeffding1948}. \citet{Sobol1993} applied this decomposition
to sensitivity analysis, showing that any square-integrable function of independent
uniform inputs admits a unique functional ANOVA into orthogonal components, and
defining the resulting variance-based indices. This decomposition gives rise to
first-order Sobol indices $S_i = D_i / D$ and total-effect indices $S_i^T$
\citep{Homma1996}, where $D$ is the total output variance and $D_i$ is the
variance contributed by input $X_i$ alone. The indices have been used extensively for factor-fixing (inputs with
$S_i^T < \tau$ can be set to any fixed value without appreciably affecting
output variance) and factor prioritisation
\citep{Saltelli2008,SaltelliStatSci2000,SaltelliInPractice2004}.

Practical estimation of Sobol indices requires $O(N(d+2))$ model evaluations
for the sample-reuse estimators of \citet{Saltelli2002} and \citet{Jansen1999}.
\citet{Saltelli2002} established the now-standard Saltelli sampling design
(also called the radial or $\mathrm{AB}$ design); \citet{Jansen1999} provided
the estimator with the best known finite-sample bias properties.
\citet{Pianosi2016} and \citet{Razavi2021} provide comprehensive reviews of
the state of the art.

The question of how transformations of the output affect sensitivity indices
has received surprisingly little systematic treatment. \citet{Saltelli2008}
note without proof that nonlinear transformations change indices, and
\citet{Puy2022} show that commonly used normalisations can lead to misleading
conclusions. \citet{Borgonovo2007} introduced the delta-index as a
transformation-invariant alternative, but it has seen limited adoption in
practice. The closest prior treatment of invariance conditions is
\citet{Borgonovo2014}, who establish qualitatively that variance-based
sensitivity statistics are invariant under affine transforms but not under
general nonlinear ones; however, they do not provide a complete biconditional
proof characterising precisely when commutativity holds.
The rank/probability-integral transform (PIT) provides a related avenue: replacing
raw outputs with their empirical quantiles before sensitivity estimation yields
rank-correlation-based SA measures \citep{Iman1979} that are invariant to
monotone transformations by construction, but as we show in
Proposition~\ref{prop:rank}, this invariance comes at the cost of commuting with
the Sobol ratio only approximately. Unified accounts of the progression from
screening to quantitative SA are given by \citet{Campolongo2011} and
\citet{PlischkeBorgonovo2013}.
The question of when and how much indices change under specific
applied transformations---the subject of this paper---remains largely unaddressed.

\subsection{Pointwise Sobol indices for spatial outputs}
\label{sec:lit_spatial}

When a model produces a spatially distributed output $Y(\mathbf{z})$, the
natural extension is to compute Sobol indices pointwise:
$S_i(\mathbf{z}) = D_i(\mathbf{z}) / D(\mathbf{z})$ for each spatial location
$\mathbf{z}$. This was developed in detail by \citet{MarrelSpatial2011} for
two-dimensional spatial benchmarks, establishing the analytical framework and
providing the Campbell2D benchmark function that has since become the standard
test case \citep{Marrel2011}. Pointwise spatial Sobol indices reveal which inputs control the model
at different locations, a capability that scalar summaries (such as the
variance-weighted aggregate of \citealt{MarrelSpatial2011}) necessarily compress.
\citet{Campbell2006} earlier treated the related problem of sensitivity analysis
when model outputs are functional, providing motivation for the pointwise approach.

Practical computation of spatial Sobol fields requires efficient
emulators~\citep{Oakley2004} or polynomial chaos
expansions~\citep{Sudret2008} when the spatial dimension is large, because
running a separate estimator for each grid cell is computationally intractable.
Sparse-grid and tensor-decomposition methods have been proposed for high-dimensional
spatial outputs \citep{Buzzard2012,Blatman2011}. A key limitation of these
approaches is that they require the sensitivity field to be smooth enough to
admit low-rank approximations; when the spatial sensitivity structure is
heterogeneous (as in our Campbell3D benchmark), more evaluations are needed.

The Campbell3D benchmark introduced in this paper extends the 2D framework
to three spatial dimensions with heterogeneous sensitivity weighting, amplifying
the spatial contrast that drives change-of-support effects.

\subsection{Sensitivity analysis for functional outputs}
\label{sec:lit_functional}

Generalising to functional outputs $Y(t)$---time series, ODE trajectories,
spectral curves---requires treating the index as a function of the independent
variable: $S_i(t) = \mathrm{Var}_{X_i}[\mathbb{E}[Y(t) \mid X_i]] / \mathrm{Var}[Y(t)]$.
The general case of sensitivity analysis for functional model outputs was treated
by \citet{Campbell2006}, who proposed aggregating pointwise variance ratios and
analysed the conditions under which a scalar summary captures global input
influence. \citet{Lamboni2011} extended this to dynamic (ODE-driven) models,
introducing a generalised sensitivity index for multivariate outputs.
The approach was formalised in the context of chemical ODE
systems by \citet{Saltelli2005}, and \citet{Sudret2008} extended it using polynomial chaos to provide
efficient estimates of $S_i(t)$ for the Damped Oscillator benchmark.
A key finding in this literature is that the sensitivity structure changes over
time: for the Damped Oscillator, stiffness $k$ dominates early in the trajectory
(when amplitude is high) while damping $\zeta$ and mass $m$ dominate at late
times (when amplitude has decayed). This time-varying sensitivity motivates the
need to carefully choose which temporal summary statistic to use for GSA.

\citet{Sudret2008} and subsequent work on polynomial chaos for functional
sensitivity \citep{Blatman2011,Crestaux2009,XiuKarniadakis2002} show that representing $S_i(t)$ without further
reduction is possible but requires $O(d \cdot N \cdot n_t)$ evaluations.
When a practitioner compresses $Y(t) \to Z = T(Y)$ (taking the peak, integrating
the trajectory, identifying threshold exceedance), they are implicitly changing
the question from ``which input most controls the full trajectory?'' to ``which
input most controls this compressed statistic?''. The answer can be dramatically
different, a key point we demonstrate empirically in Section~\ref{sec:temporal_results}.

\subsection{Transformations, commutativity, and the current gap}
\label{sec:lit_gap}

The relationship between output transformations and sensitivity indices has been
studied in specific contexts. \citet{Pianosi2016} discusses the connection between
sensitivity indices and the choice of output metric, noting that different
metrics can lead to different rankings but without systematic quantification.
\citet{Puy2022} show that rescaling outputs by their mean can reverse sensitivity
rankings, demonstrating a concrete case of noncommutativity. \citet{Oakley2004}
discuss how Gaussian process emulators can represent the full output distribution
and thus support sensitivity analysis under arbitrary output transformations,
but leave the commutativity question implicit.

More fundamentally, \citet{Saltelli2008} and \citet{IoossLemaitre2015} note
that the choice of what to analyse (the model output or a function of it) is
a decision that should be made explicitly and consistently. Yet applied papers
routinely report indices for a transformed output without acknowledging or
accounting for the transformation.

The framework we introduce in Section~\ref{sec:framework}---FSAO---provides
a formal language for this decision, a rigorous commutativity criterion
(Propositions~\ref{prop:main}--\ref{prop:scalar}), and the first comprehensive
empirical quantification of the consequences of the decision across multiple
output types and transform categories.

%% ─────────────────────────────────────────────────────────────────────────────
\section{Framework: Functional Sensitivity Analysis Under Output Operators}
\label{sec:framework}
%% ─────────────────────────────────────────────────────────────────────────────

\subsection{Unified setup and nomenclature}
\label{sec:nomenclature}

Let $f: \Omega \to \mathcal{Y}$ where $\Omega = \prod_{i=1}^d [a_i, b_i]$ and
$\mathcal{Y}$ is a Hilbert space of real-valued functions on an index set
$\mathcal{Z}$.  The index set $\mathcal{Z}$ is:
\begin{itemize}[noitemsep]
  \item A finite set $\{z_1, \ldots, z_p\}$, yielding a scalar output ($p=1$)
    or a finite-dimensional vector output ($p > 1$).
  \item A spatial grid $\mathcal{Z} \subset \mathbb{R}^2$ or $\mathbb{R}^3$,
    yielding a \emph{spatial functional output}.
  \item A time grid $\{t_0, t_1, \ldots, t_{n_t-1}\} \subset \mathbb{R}_{\geq 0}$,
    yielding a \emph{temporal functional output}.
\end{itemize}
We write $Y_z = f(\mathbf{X})(z)$ for the output at index point $z \in \mathcal{Z}$.
With inputs $\mathbf{X} \sim \mathrm{Uniform}(\Omega)$, the first-order Sobol
index at $z$ is
\begin{equation}
  S_i(z) = \frac{\mathrm{Var}_{X_i}[\mathbb{E}_{\mathbf{X}_{\sim i}}[Y_z \mid X_i]]}
                 {\mathrm{Var}[Y_z]},
  \label{eq:pointwise_sobol}
\end{equation}
and the scalar benchmark ($|\mathcal{Z}|=1$) is a special case.

\begin{definition}[Output operator]
An \emph{output operator} $T: \mathcal{Y} \to \mathcal{Y}'$ is any map that
takes the functional output $Y \in \mathcal{Y}$ and produces a new functional
output $Z = T(Y) \in \mathcal{Y}'$.  We say $T$ is:
\begin{itemize}[noitemsep]
  \item \emph{pointwise} if $T(Y)(z) = \phi(Y(z))$ for a fixed function
    $\phi: \mathbb{R} \to \mathbb{R}$;
  \item \emph{aggregating} if $|\mathcal{Z}'| < |\mathcal{Z}|$
    (e.g.\ block averaging, temporal peak, exceedance area);
  \item \emph{differentiating} if it involves derivatives or differences of
    $Y$ across $\mathcal{Z}$ (e.g.\ gradient magnitude, Laplacian roughness);
  \item \emph{spectral} if it operates in the frequency/transform domain
    (e.g.\ bandpass filter).
\end{itemize}
\end{definition}

\begin{definition}[Commutativity]
\label{def:commute}
The operator $T$ \emph{commutes with Sobol index computation} if
$S_i^{T \circ f}(z') = S_i^f(z)$ for all $i \in \{1,\ldots,d\}$ and all
$z' \in \mathcal{Z}'$, where $z'$ is the image of $z$ under $T$.
When $|\mathcal{Z}'| < |\mathcal{Z}|$, commutativity is understood in the
sense of a variance-weighted aggregation:
\begin{equation}
  \hat{S}_i(T \circ f) = \frac{\sum_{z' \in \mathcal{Z}'} \mathrm{Var}[Z_{z'}]\, S_i^{T\circ f}(z')}
                               {\sum_{z'} \mathrm{Var}[Z_{z'}]}
  = \hat{S}_i(f) := \frac{\sum_{z \in \mathcal{Z}} \mathrm{Var}[Y_z]\, S_i^f(z)}
                         {\sum_z \mathrm{Var}[Y_z]}.
  \label{eq:commute_def}
\end{equation}
\end{definition}

The question ``does $T$ commute with Sobol?'' is thus a well-posed mathematical
statement for any output type, not just scalars.
We use ``non-commutativity'' to refer to the failure of \eqref{eq:commute_def}.

\subsection{The commutativity theorem}
\label{sec:commute_thm}

\begin{proposition}[Commutativity iff affine]
\label{prop:main}
Let $T$ be a pointwise operator $T(Y)(z) = \phi(Y_z)$ for a fixed function
$\phi: \mathbb{R} \to \mathbb{R}$.  Then $T$ commutes with Sobol index
computation in the sense of Definition~\ref{def:commute} if and only if
$\phi$ is affine, i.e.\ $\phi(y) = ay + b$ for constants $a \neq 0$, $b$.
\end{proposition}

A complete proof is given in Appendix~\ref{app:proofs}.  The key idea is that
$Z_z = aY_z + b$ gives $\mathrm{Var}[Z_z] = a^2 \mathrm{Var}[Y_z]$ and
$\mathrm{Var}_{X_i}[\mathbb{E}[Z_z \mid X_i]] = a^2 \mathrm{Var}_{X_i}[\mathbb{E}[Y_z \mid X_i]]$,
so the ratio $S_i(z)$ is unchanged.  Conversely, Jensen's inequality shows
that a strictly convex $\phi$ gives $\mathbb{E}[\phi(Y_z) \mid X_i] \neq \phi(\mathbb{E}[Y_z \mid X_i])$,
breaking the variance ratio.

\begin{remark}[Exact scalar counterexample]
\label{rem:toy}
The non-commutativity mechanism is transparent in a minimal example requiring
no Monte Carlo.  Let $Y = X_1 + X_2$ with $X_1, X_2 \stackrel{iid}{\sim}
\mathrm{Bernoulli}(1/2)$, so $S_1 = S_2 = 1/2$ exactly (additive model,
equal variances).  Now apply $\phi(y) = y^2$.  Then
$Z = (X_1 + X_2)^2 = X_1^2 + 2X_1 X_2 + X_2^2 = X_1 + 2X_1 X_2 + X_2$
(using $X_i^2 = X_i$ for Bernoulli inputs).  The transformed model is no
longer additive: the interaction term $2X_1 X_2$ appears, redistributing
variance.  Explicitly, $\mathrm{Var}[Z] = 3/4$,
$D_1^Z = D_2^Z = 1/4$, and $S_1^Z = S_2^Z = 1/3 \neq 1/2$.  The
total-effect indices shift correspondingly: $S_1^{T,Z} = S_2^{T,Z} = 2/3$.
This illustrates the core mechanism: a nonlinear $\phi$ creates interaction
structure absent from the original model, changing both first-order and
total-effect decompositions.  The full proof in Appendix~\ref{app:proofs}
generalises this observation to arbitrary non-affine $\phi$.
\end{remark}

Propositions~\ref{prop:cos} and \ref{prop:scalar} refine this result for two
practically important special structures---the change-of-support operator and scalar
outputs under monotone transforms---and are stated within the Results section
(Sections~\ref{sec:c3d} and~\ref{sec:scalar_results}) where the empirical evidence
that motivates them is presented.

\subsection{Variance-weighted aggregation}
\label{sec:vw_agg}

For any functional output---spatial, temporal, or otherwise---we use the
variance-weighted aggregate
\begin{equation}
  \hat{S}_i = \frac{\sum_{z \in \mathcal{Z}} \mathrm{Var}[Y_z]\, \hat{S}_i(z)}
                   {\sum_{z \in \mathcal{Z}} \mathrm{Var}[Y_z]}
  \label{eq:vw}
\end{equation}
as the scalar summary for comparing the original and transformed index vectors.
This weighting was proposed by \citet{MarrelSpatial2011} for spatial outputs and
extends naturally to temporal outputs: $z$ is a time point $t$ rather than a
spatial coordinate, and $\mathrm{Var}[Y_t]$ is the cross-sample variance of
the trajectory at time $t$.

The interpretation is consistent across output types: $\hat{S}_i$ measures
the fraction of the total functional output variance attributable to $X_i$
at the average location (weighted by local variance contribution).
For scalar outputs, $\hat{S}_i = S_i$ exactly.
For functional outputs, $\hat{S}_i$ is a weighted mean over the index set,
placing more weight on high-variance index points where sensitivity is most
consequential.

We emphasise that Proposition~\ref{prop:main} is fundamentally a
\emph{pointwise} invariance result: $S_i(z)$ is unchanged by an affine $\phi$
at every $z$ independently.  The variance-weighted aggregate enters only as a
\emph{reporting convention} for summarising the pointwise field; it does not
appear in the theorem's hypotheses. Because \eqref{eq:vw} is a convex
combination of pointwise indices, pointwise commutativity implies aggregate
commutativity, but the converse does not hold in general.

%% ─────────────────────────────────────────────────────────────────────────────
\section{Experimental Setup}
\label{sec:setup}
%% ─────────────────────────────────────────────────────────────────────────────

\subsection{Benchmark functions}
\label{sec:benchmarks}

All benchmarks are implemented in the \textit{sabench} library~\citep{Hettinger2025}.
We group them by output type (Table~\ref{tab:benchmarks}).

\begin{table}[H]
\centering
\caption{Benchmark functions. $d$ = input dimension; $N_\text{ref}$ = refinement
sample size; ``type'' = scalar / spatial / temporal.}
\label{tab:benchmarks}
\small
\begin{tabular}{llccp{5cm}}
\toprule
Benchmark & Type & $d$ & $N_\text{ref}$ & Key features \\
\midrule
\multicolumn{5}{l}{\textit{Spatial benchmarks}} \\
Campbell2D & Spatial $32^2$ & 8 & 2048 & Heterogeneous $S_i(\mathbf{z})$;
  $V_5\equiv0$; $V_6=V_8$ \\
Campbell3D & Spatial $12^3$ & 8 & 1024 & 3D additive mixing; analytical $S_i$ available; $V_5\equiv0$; $V_8=V_6$ \\
\midrule
\multicolumn{5}{l}{\textit{Temporal / functional benchmarks}} \\
DampedOscillator & Temporal $n_t=100$ & 3 & 1024 & $S_k(t)$ dominant early;
  $S_\zeta(t)$ dominant late \\
LotkaVolterra    & Temporal $n_t=100$ & 4 & 1024 & ODE; cyclic sensitivity \\
\midrule
\multicolumn{5}{l}{\textit{Scalar benchmarks}} \\
Ishigami   & Scalar & 3 & 2048 & Strong interaction; $X_3$ inactive \\
SobolG     & Scalar & 8 & 2048 & Many near-threshold inputs; high vulnerability \\
Borehole   & Scalar & 8 & 2048 & 2 dominant inputs; 6 near-zero \\
Piston     & Scalar & 7 & 2048 & Moderate interactions \\
WingWeight & Scalar & 10 & 2048 & Near-additive; 10 inputs \\
OTL Circuit & Scalar & 6 & 2048 & Electronic; mixed interactions \\
\bottomrule
\end{tabular}
\end{table}

\paragraph{Spatial benchmarks.}
The Campbell2D benchmark \citep{MarrelSpatial2011,Marrel2011} is a $32\times32$ spatial field
driven by 8 inputs, with a known heterogeneous sensitivity structure.
Two structural invariants serve as verification: $V_5 \equiv 0$ (input $X_5$
has zero first-order effect everywhere) and $V_6 = V_8$ (inputs $X_6$ and $X_8$
are interchangeable by symmetry).

The Campbell3D benchmark is an extension we introduce here for the first time
\citep{Hettinger2025c3d}.  It generalises Campbell2D to a third spatial dimension
by appending an independent additive mixing term $c_{i3}z_3$ to each 2D mixing
coordinate, preserving exact backward compatibility ($g_3(\mathbf{X},z_1,z_2,0)
\equiv g_2(\mathbf{X},z_1,z_2)$) while avoiding the heterogeneity collapse that
afflicts the natural sum-to-one extension.  The output is a $12\times12\times12$
volume for 8 inputs.  Closed-form first-order partial variances are available
via the same Gaussian-quadrature structure as Campbell2D
\citep{MarrelSpatial2011}---the 3D mixing coordinates enter the partial-variance
integrals as arguments only, so the structural invariants
$V_5 \equiv 0$ and $V_8 = V_6$ are preserved exactly.  The benchmark has not
previously appeared in the sensitivity analysis literature to our knowledge.

\paragraph{Temporal benchmarks.}
The Damped Oscillator \citep{Sudret2008} produces a displacement trajectory
$x(t) = (F/m\omega_d)\exp(-\zeta\omega_n t)\sin(\omega_d t + \phi)$, discretised at
$n_t = 100$ time points over $[0, 10]$\,s, for 3 inputs
$(m, k, \zeta)$.  Its pointwise sensitivity structure has a known qualitative
property: stiffness $k$ dominates the early trajectory (when the exponential
envelope is large), while damping $\zeta$ and mass $m$ take over at later times.

The Lotka-Volterra predator-prey ODE \citep{Lotka1925,Volterra1926}, used here
as a sensitivity benchmark following \citet{Saltelli2005}, produces the prey
population trajectory $N(t)$ for 4 inputs $(\alpha, \beta, \delta, \gamma)$,
solved by RK4 over $[0, 20]$.  Its sensitivity structure is cyclic: the prey
growth rate $\alpha$ dominates near population maxima while predation efficiency
$\delta$ matters near minima.

\paragraph{Scalar benchmarks.}
The Ishigami function \citep{Ishigami1990} is the canonical nonlinear SA test case
with known analytical indices. The SobolG function was introduced by
\citet{SobolLevitan1999} to stress-test estimators on near-threshold inputs.
The Borehole model originates from contaminant transport simulation; its sensitivity
structure was characterised by \citet{MorrisMitchell1993}. The Piston and OTL
Circuit functions are engineering models from \citet{BenAriSteinberg2007}. The
WingWeight function is the 10-input aeronautics model from \citet{Forrester2008}.
All six have well-established first-order Sobol indices \citep{Saltelli2008}, and
we use them as a reference group for the scalar invariance theorem
(Section~\ref{sec:scalar_results}).

\subsection{Transform registry}
\label{sec:transforms}

We register 33 transforms in 5 categories (Table~\ref{tab:transforms}).
The temporal category is new in this work.

\begin{table}[H]
\centering
\caption{Transform registry. Cat: Env=environmental, Eng=engineering, Spa=spatial,
Tmp=temporal, Sta=statistical. $D_\mu$ and $\Delta_\mu$: cross-benchmark means.
${}^\dagger$: linear transform (Proposition~\ref{prop:main} predicts $D \approx 0$).}
\label{tab:transforms}
\footnotesize
\begin{tabular}{llccc}
\toprule
Key & Name & Cat. & $D_\mu$ & $\Delta_\mu$ \\
\midrule
\texttt{arrhenius}                     & Arrhenius $\exp(-E_a/Ry)$ & Eng & 0.372 & 0.606 \\
\texttt{exceedance\_area}              & Exceedance area fraction & Spa & 0.405 & 0.748 \\
\texttt{isoline\_length}               & Isoline perimeter & Spa & 0.405 & 0.669 \\
\texttt{log\_shift}                    & Log-shift $\log(y+c)$ & Env & 0.364 & 0.594 \\
\texttt{box\_cox\_log}                 & Box-Cox $\lambda\to0$ & Env & 0.364 & 0.594 \\
\texttt{gradient\_magnitude}           & Gradient magnitude $|\nabla Y|$ & Spa & 0.391 & 0.648 \\
\texttt{laplacian\_roughness}          & Laplacian $|\nabla^2 Y|$ & Spa & 0.388 & 0.660 \\
\texttt{box\_cox\_sqrt}                & Box-Cox $\lambda=0.5$ & Env & 0.349 & 0.571 \\
\texttt{clipped\_excess\_q90}          & Clipped excess POT & Env & 0.367 & 0.626 \\
\texttt{power\_law\_beta2}             & Power-law $y^2$ & Env & 0.378 & 0.612 \\
\texttt{sigmoid\_dose}                 & Hill dose-response & Eng & 0.378 & 0.630 \\
\texttt{normalised\_stress}            & Goodman normalised stress & Eng & 0.359 & 0.608 \\
\texttt{power\_law\_beta05}            & Power-law $y^{0.5}$ & Env & 0.370 & 0.620 \\
\texttt{contour\_exceedance}           & Signed contour distance & Spa & 0.372 & 0.610 \\
\texttt{carnot\_quadratic}             & Carnot quadratic & Eng & 0.369 & 0.608 \\
\texttt{weibull\_reliability}          & Weibull failure prob. & Eng & 0.378 & 0.617 \\
\texttt{exceedance\_q75}               & Exceedance $q_{75}$ & Env & 0.363 & 0.630 \\
\texttt{exceedance\_q90}               & Exceedance $q_{90}$ & Env & 0.366 & 0.615 \\
\texttt{exceedance\_q95}               & Exceedance $q_{95}$ & Env & 0.369 & 0.620 \\
\texttt{exceedance\_q99}               & Exceedance $q_{99}$ & Env & 0.360 & 0.624 \\
\texttt{entropy\_proxy}                & Entropy proxy $-(Y-\bar Y)^2$ & Sta & 0.337 & 0.558 \\
\texttt{standardised\_anomaly}         & Standardised anomaly & Sta & 0.364 & 0.609 \\
\texttt{temporal\_log\_cumsum}         & Log cumulative sum & Tmp & 0.189 & 0.338 \\
\texttt{temporal\_exceedance\_duration}& Exceedance duration & Tmp & 0.219 & 0.435 \\
\texttt{temporal\_peak}                & Temporal peak $\max_t$ & Tmp & 0.204 & 0.454 \\
\texttt{temporal\_envelope}            & Running maximum envelope & Tmp & 0.185 & 0.427 \\
\texttt{temporal\_bandpass}${}^\dagger$& Bandpass filter (linear) & Tmp & 0.301 & 0.527 \\
\texttt{temporal\_block\_avg}          & Temporal block avg (COS) & Tmp & 0.019 & 0.069 \\
\texttt{temporal\_cumsum}${}^\dagger$  & Cumulative sum (linear) & Tmp & 0.146 & 0.479 \\
\texttt{regional\_mean}                & Spatial regional mean & Spa & 0.048 & 0.128 \\
\texttt{block\_8x8}                    & Block average $8^3$ & Spa & 0.013 & 0.028 \\
\texttt{block\_4x4}                    & Block average $4^3$ & Spa & 0.003 & 0.007 \\
\texttt{block\_2x2}                    & Block average $2^3$ & Spa & 0.001 & 0.002 \\
\texttt{matern\_smooth}                & Matérn-$\frac{5}{2}$ smooth & Spa & 0.002 & 0.005 \\
\texttt{rank\_transform}               & Rank / PIT & Sta & 0.042 & 0.104 \\
\bottomrule
\end{tabular}
\end{table}

\paragraph{Transform provenance.}
Each engineering and environmental transform is grounded in an established
physical or statistical model.
The Arrhenius transform $\exp(-E_a/Ry)$ originates in chemical kinetics
\citep{Arrhenius1889}.
The Hill sigmoid dose-response $y^n/(k^n + y^n)$ follows \citet{Hill1910}.
The Weibull failure probability $1 - \exp(-(y/\lambda)^k)$ follows
\citet{Weibull1951}.
The Goodman normalised stress transform is the fatigue-life criterion from
structural engineering \citep{SocieMarquis2000}.
The Box-Cox family $\lambda \in \{0, 0.5\}$ follows \citet{BoxCox1964}.
The Matérn-$\tfrac{5}{2}$ spatial smoother applies a kernel derived from the
Matérn covariance class \citep{Matern1986}.
Log-shift and Box-Cox transforms are routinely applied to hydrological model
outputs (e.g.\ streamflow) before calibration and GSA; a recent systematic
analysis of their effect on model performance is given by \citet{Thirel2024}.
The rank/PIT statistical transform replaces raw values with empirical quantiles
following \citet{Iman1979}.

\paragraph{Rationale for Arrhenius as reference transform.}
Several sections use the Arrhenius transform $T(y) = \exp(-E_a/R \cdot y_\text{pos})$
as the primary illustration, where the positivity shift
$y_\text{pos} = y - y_\text{min} + \varepsilon$ uses a \emph{fixed} constant
$y_\text{min}$ determined from the \emph{analytic support} of each benchmark
(not from the sample), ensuring that $T$ remains a fixed pointwise map
$\phi: \mathbb{R} \to \mathbb{R}$ as required by Proposition~\ref{prop:main}.
This choice is deliberate: (i) it is parametric
and physically motivated (atmospheric chemistry, thermal diffusivity, materials
science \citep{Arrhenius1889}), (ii) its degree of nonlinearity is controlled by
$E_a/R$, making it a tunable reference, (iii) it consistently achieves among the
highest $D$ and $\Delta$ across all benchmarks (mean $D = 0.372$), making it a
useful worst-case indicator. We do not suggest that Arrhenius transformations are uniquely important;
log-shift, Box-Cox \citep{BoxCox1964}, and power-law transforms produce nearly
identical scores on scalar benchmarks (scalar invariance theorem), and the
choice is purely for consistency of illustration.

\subsection{Scoring metrics}
\label{sec:metrics}

We score each (benchmark, transform) pair using two bounded metrics.
Both were designed to overcome specific failure modes of simpler alternatives.
Preliminary experiments used a composite score combining integer threshold-flip
counts and a relative-$L^2$ norm; this failed in two ways: the flip counts
saturated at their maximum for 20 of 26 transforms (degeneracy, producing
identical composite values near 14.7), and the relative-$L^2$ diverged to
$O(10^{11})$ on the Borehole benchmark where 6 of 8 inputs have $\hat{S}_i \approx 0$.
These lessons motivated the two metrics below.

\paragraph{Decision Score $D$.}
The Decision Score measures the classification impact of the transform on
factor-fixing decisions. Define the soft sigmoid threshold
\begin{equation}
  \sigma(s; \tau) = \frac{1}{1 + \exp(-(s - \tau)/\tau)},
  \label{eq:soft}
\end{equation}
which transitions from 0 (inactive) to 1 (active) around the practical
significance threshold $\tau = 0.05$ \citep{Saltelli2008}.
The Decision Score is the mean per-input absolute difference in soft-classification:
\begin{equation}
  D = \frac{1}{d}\sum_{i=1}^d \bigl|\sigma(\hat{S}_i(Z); \tau) - \sigma(\hat{S}_i(Y); \tau)\bigr|.
  \label{eq:D}
\end{equation}
$D = 0$ means no input's active/inactive classification changed; $D = 1$ would
require every input to flip from deep inactive to deep active simultaneously.
Intuitively, $D$ asks whether any input crossed the practical significance
threshold $\tau$ as a result of the transform, and averages those soft
threshold-crossings over all inputs.
The sigmoid form avoids the degeneracy of integer flip counts while preserving
the decision-threshold interpretation. The threshold $\tau = 0.05$ and the
factor-fixing framing follow \citet{Saltelli2008}; the soft-classification idea
is inspired by the screening literature, where inputs are classified as active
or inactive relative to a fixed threshold \citep{Morris1991,Campolongo2007}.
The sigmoid width equal to $\tau$ provides natural calibration around the
boundary. Unlike a rank-correlation measure (e.g., Spearman $\rho$), $D$ is
sensitive to absolute index changes even when rankings are preserved, which
matters when inputs cross the significance threshold.

\paragraph{Sensitivity Shift $\Delta$.}
The Sensitivity Shift measures raw redistribution of sensitivity mass, using the
Bray-Curtis dissimilarity \citep{Bray1957,Ferretti2016}:
\begin{equation}
  \Delta = \frac{\sum_i |\hat{S}_i(Z) - \hat{S}_i(Y)|}
                {\sum_i \hat{S}_i(Z) + \sum_i \hat{S}_i(Y)}.
  \label{eq:delta}
\end{equation}
The pooled denominator makes $\Delta$ robust when $\hat{S}_i(Y) \approx 0$
for many inputs, which caused the relative-$L^2$ divergence on Borehole.
Bray-Curtis dissimilarity is the standard measure of community composition
change in ecology \citep{Bray1957} and has been applied to sensitivity index
vectors \citep{Ferretti2016}; its bounded $[0,1]$ range makes it
directly comparable across benchmarks and transform categories.
Intuitively, $\Delta$ asks what fraction of the total sensitivity mass had to
move to convert the original index profile into the transformed one: $\Delta = 0$
means the transform left the profile entirely intact and $\Delta = 1$ means the
two profiles share no overlap.

\paragraph{Estimation.}
For each benchmark $f$ and transform $T$, we generate a Saltelli sample
$X \in \mathbb{R}^{N(d+2) \times d}$ \citep{Saltelli2002}, evaluate $Y = f(X)$
and $Z = T(Y)$, apply the Jansen (1999) estimator~\citep{Jansen1999} to both,
and compute $\hat{S}_i$ via variance weighting (Eq.~\ref{eq:vw}).
A two-stage screening/refinement design at $N=512$ and $N=2048$
(or $N=1024$ for memory-heavy benchmarks) manages computational cost.

\paragraph{First-order versus total-effect indices.}
Throughout this paper we report first-order Sobol indices $S_i$.
Total-effect indices $S_i^T$ \citep{Homma1996} quantify interaction-inclusive
sensitivity and are often preferred for factor fixing.  We focus on $S_i$
because (i)~they are estimable from the same Saltelli sample at no additional
cost, (ii)~many environmental benchmarks are near-additive so $S_i \approx S_i^T$,
and (iii)~the non-commutativity argument is cleanest for the additive component.
The analogous results for total-effect indices---including numerical confirmation
that $D$ and $\Delta$ are elevated for the same transforms---are presented in
Appendix~\ref{app:total_effects}; main-text conclusions are unchanged.

%% ─────────────────────────────────────────────────────────────────────────────
\section{Results}
\label{sec:results}
%% ─────────────────────────────────────────────────────────────────────────────

\subsection{Overview: heatmaps across all benchmarks}
\label{sec:heatmaps}

Figure~\ref{fig:heatmaps} shows the $D$ and $\Delta$ heatmaps for all
33 transforms $\times$ 10 benchmarks. The dominant structural pattern is
a clean separation between the top 20 transforms (environmental and engineering
nonlinear transforms: $D \in [0.2, 0.6]$) and the bottom group (spatial COS
operators, linear temporal transforms: $D < 0.1$). Temporal transforms occupy
an intermediate position.

\begin{figure}[ht]
\centering
\begin{subfigure}[t]{0.48\textwidth}
  \includegraphics[width=\textwidth]{fig_heatmap_D_s1.png}
  \caption{Decision Score $D$. Grey = skipped (C3D Matérn).}
\end{subfigure}
\hfill
\begin{subfigure}[t]{0.48\textwidth}
  \includegraphics[width=\textwidth]{fig_heatmap_delta_s1.png}
  \caption{Sensitivity Shift $\Delta$.}
\end{subfigure}
\caption{Non-commutativity heatmaps for 33 transforms $\times$ 10 benchmarks.
  Transforms sorted by mean $D$ (top = most noncommutative). Y-axis label colours:
  {\color{cat_env}blue}=environmental, {\color{cat_eng}orange}=engineering,
  {\color{cat_spa}green}=spatial, {\color{cat_tmp}purple}=temporal,
  {\color{cat_sta}grey}=statistical.
  X-axis label colours indicate benchmark type (green=spatial, purple=temporal, grey=scalar).}
\label{fig:heatmaps}
\end{figure}

\subsection{Spatial results}
\label{sec:spatial_results}

\subsubsection{Campbell2D: per-input structure and spatial maps}

Figure~\ref{fig:per_input_spatial} shows per-input $\hat{S}_i$ for Campbell2D
under the original output and three selected transforms.

\begin{figure}[ht]
\centering
\includegraphics[width=0.9\textwidth]{fig_per_input_spatial.png}
\caption{Per-input $\hat{S}_i$ on Campbell2D ($N=2048$) for original output and
  three transforms. Structural invariants $V_5 \equiv 0$ and $V_6 = V_8$ hold
  exactly in the analytical model; the Monte Carlo estimates respect them to
  within estimator noise ($|\hat{S}_5| < 0.005$; $|\hat{S}_6 - \hat{S}_8| < 0.04$
  at $N=2048$). Inputs 1, 2, 4 show the largest shifts under
  Arrhenius and log-shift; block-$8\times8$ shows milder but nontrivial shifts.}
\label{fig:per_input_spatial}
\end{figure}

The spatial Sobol maps in Figure~\ref{fig:spatial_maps} illustrate the pointwise
changes.  On the original output, sensitivity is distributed across $X_2$,
$X_4$, $X_6$, and $X_8$, while $X_1$ is comparatively weak throughout the
field.  The Arrhenius exponential is most sensitive to variation in $Y$ near
the lower end of the output range, precisely where $X_1$ exerts its influence;
the result is a sharp amplification of $S_1(\mathbf{z})$ at those locations
and a simultaneous collapse of $S_2$, $S_4$, $S_6$, and $S_8$ relative to the
post-transform total variance.

\begin{figure}[ht]
\centering
\includegraphics[width=\textwidth]{fig_spatial_maps_v3.png}
\caption{Pointwise Sobol index maps $S_i(\mathbf{z})$ for Campbell2D, inputs
  $X_1$--$X_8$. Original output (top row) vs Arrhenius transform (bottom row).
  On the original output $X_2$, $X_4$, $X_6$, and $X_8$ show diffuse
  distributed sensitivity while $X_1$ is weak throughout; after Arrhenius,
  $X_1$ becomes the single dominant input and all four originally dominant
  inputs collapse to near zero.
  The analytical structural invariants $S_5 \equiv 0$ and $S_6 = S_8$ are
  respected to within Monte Carlo noise in both rows.}
\label{fig:spatial_maps}
\end{figure}

Figure~\ref{fig:contour_maps} shows the new contour-exceedance transform:
a signed distance field relative to the $q_{75}$ contour level, and its
effect on the Sobol field.

\begin{figure}[ht]
\centering
\includegraphics[width=\textwidth]{fig_contour_maps.png}
\caption{Contour exceedance transform on Campbell2D. Left: an example output
  field with the $q_{75}$ contour overlay (red). Centre-left: the transformed
  field (signed distance, blue=below, red=above contour). Centre-right and right:
  $S_1(\mathbf{z})$ under original vs contour-exceedance. The transform
  concentrates sensitivity in the contour transition zone.}
\label{fig:contour_maps}
\end{figure}

\subsubsection{Campbell3D: novel benchmark, stronger COS effects}
\label{sec:c3d}

The Campbell3D benchmark introduces a third spatial mixing dimension with
heterogeneous weighting, amplifying sensitivity contrasts across the volume.
This is, to our knowledge, the first published sensitivity analysis of this
benchmark. It was developed as part of \textit{sabench} and has not appeared
in the prior literature.

Figure~\ref{fig:c2d_c3d} compares the Decision Score on Campbell2D and
Campbell3D for the top-20 transforms. Two patterns are notable:

\begin{figure}[ht]
\centering
\includegraphics[width=0.85\textwidth]{fig_c2d_vs_c3d.png}
\caption{Campbell2D vs Campbell3D Decision Score $D$ for top-20 transforms.
  Environmental/engineering transforms score slightly lower on C3D due to its
  more diffuse index profile. Spatial block-averaging (COS) scores dramatically
  higher on C3D: block-$8^3$ gives $D = 0.079$ on C3D vs $D = 0.006$ on C2D.}
\label{fig:c2d_c3d}
\end{figure}

\textit{Environmental/engineering transforms score somewhat lower on C3D.}
The C3D benchmark has a more uniform $\hat{S}_i$ profile across the volume
(the 3D mixing spreads sensitivity more evenly), meaning fewer inputs are
near the $\tau=0.05$ threshold simultaneously.  With fewer near-threshold
inputs available to flip, the Decision Score is smaller.

\textit{Spatial COS transforms score dramatically higher on C3D.}
The spatial heterogeneity of $S_i(\mathbf{z})$ is amplified by the third
dimension, so block averaging redistributes sensitivity mass much more strongly.
Table~\ref{tab:cos} gives the quantitative comparison.

\begin{table}[H]
\centering
\caption{Change-of-support scaling: $D$ and $\Delta$ by block scale.
Campbell3D at $8^3$ is $13.2\times$ larger than Campbell2D at $8\times8$.}
\label{tab:cos}
\begin{tabular}{lcccc}
\toprule
& \multicolumn{2}{c}{Campbell2D} & \multicolumn{2}{c}{Campbell3D} \\
\cmidrule(lr){2-3}\cmidrule(lr){4-5}
Block & $D$ & $\Delta$ & $D$ & $\Delta$ \\
\midrule
$2\times2$ ($2^3$) & 0.0003 & 0.0008 & 0.0031 & 0.0070 \\
$4\times4$ ($4^3$) & 0.0016 & 0.0037 & 0.0114 & 0.0283 \\
$8\times8$ ($8^3$) & 0.0060 & 0.0140 & 0.0790 & 0.1360 \\
\bottomrule
\end{tabular}
\end{table}

The data in Table~\ref{tab:cos} call for a precise theoretical characterisation
of when block-averaging commutes with Sobol index computation.

\begin{proposition}[COS commutativity condition]
\label{prop:cos}
Let $T^{\mathrm{COS}}_B(Y) = |B|^{-1}\int_B Y(z)\,dz$ be the block-average
operator over region $B \subset \mathcal{Z}$ (the change-of-support operator
of \citealt{Cressie1993}).  Then $T^{\mathrm{COS}}_B$
commutes with Sobol index computation if and only if the pointwise sensitivity
structure $S_i(z)$ is constant over $B$ for all inputs $i$,
i.e.\ $S_i(z) = S_i$ for all $z \in B$.
A complete proof is given in Appendix~\ref{app:proofs}.
\end{proposition}

The $13.2\times$ amplification from C2D to C3D at $8^3$ blocks is direct
empirical validation of Proposition~\ref{prop:cos}: Campbell3D's stronger 3D
heterogeneity means that a block of side $8$ units spans multiple distinct
sensitivity regimes, violating the homogeneity condition at a scale where
Campbell2D does not.  Scalar benchmarks have $|\mathcal{Z}|=1$, so the
sensitivity structure is trivially constant and COS scores $D=0$.
The Damped Oscillator's temporal COS at block-10 gives $D = 0.013$, reflecting
mild heterogeneity in $S_i(t)$ over short intervals.
The spatial variogram of $S_i(\mathbf{z})$ in C3D has a shorter correlation
length than in C2D (the third dimension introduces finer-scale heterogeneity),
so $8^3$ blocks span multiple distinct sensitivity regimes, producing a larger
COS effect.  The scaling curve (Figure~\ref{fig:cos_scaling}) is monotone
and smooth for both benchmarks, consistent with the theory.

\begin{figure}[ht]
\centering
\includegraphics[width=0.7\textwidth]{fig_cos_scaling.png}
\caption{COS scaling: $D$ and $\Delta$ vs block scale for Campbell2D and C3D.
  Monotone increase confirms the theory; C3D amplifies the effect.}
\label{fig:cos_scaling}
\end{figure}

\subsection{Temporal results}
\label{sec:temporal_results}

The temporal benchmarks introduce a new dimension to the non-commutativity
analysis.  For a time-series output $Y(t)$, the pointwise Sobol indices
$S_i(t)$ carry information about the time-varying dominance of each input.
Any temporal post-processing operator $T$ that compresses or reshapes this
trajectory potentially changes which inputs appear dominant in the aggregate.

\subsubsection{Damped Oscillator: linear vs nonlinear temporal transforms}

Figure~\ref{fig:temporal_traces} shows the pointwise indices $S_i(t)$ for the
Damped Oscillator under the original trajectory and three temporal transforms.

\begin{figure}[ht]
\centering
\includegraphics[width=0.95\textwidth]{fig_temporal_traces.png}
\caption{Pointwise Sobol indices $S_i(t)$ for the Damped Oscillator ($N=1024$).
  Top-left: original trajectory dominated by $k$ (stiffness) early and $\zeta$
  (damping) late. Top-right: temporal peak ($\max_t$) collapses to a single value,
  strongly favouring $k$. Bottom-left: log-cumsum redistributes toward $\zeta$
  (late-time contribution dominates the integrated log). Bottom-right: exceedance
  duration amplifies the early-time contribution of $k$.}
\label{fig:temporal_traces}
\end{figure}

\begin{figure}[ht]
\centering
\includegraphics[width=0.95\textwidth]{fig_per_input_temporal.png}
\caption{Variance-weighted aggregate $\hat{S}_i$ for the Damped Oscillator
  under four temporal transforms plus the linear cumsum control. The cumsum
  (linear) closely tracks the original, consistent with
  Corollary~\ref{cor:linear}. All nonlinear transforms produce substantially
  different index profiles.}
\label{fig:per_input_temporal}
\end{figure}

\begin{figure}[ht]
\centering
\includegraphics[width=0.95\textwidth]{fig_linear_controls.png}
\caption{Linear vs nonlinear temporal transforms: mean $D$ (solid) and $\Delta$
  (hatched) on Damped Oscillator and Lotka-Volterra. Linear transforms
  (cumsum, bandpass) score lower than nonlinear transforms, consistent with
  Proposition~\ref{prop:main}. The bandpass scores higher than expected for a
  linear transform because it selects a frequency band that differs from the
  dominant spectral content, shifting variance weights substantially.}
\label{fig:linear_controls}
\end{figure}

Table~\ref{tab:temporal} gives the numerical scores for temporal transforms
on both benchmarks.

\begin{table}[H]
\centering
\caption{Temporal transform scores on functional benchmarks. ${}^\dagger$: linear.
Bandpass is linear but scores moderately high due to variance-weight reallocation
when most energy is removed from the passband.}
\label{tab:temporal}
\small
\begin{tabular}{lcccc}
\toprule
& \multicolumn{2}{c}{DampedOscillator} & \multicolumn{2}{c}{LotkaVolterra} \\
\cmidrule(lr){2-3}\cmidrule(lr){4-5}
Transform & $D$ & $\Delta$ & $D$ & $\Delta$ \\
\midrule
Temporal peak ($\max_t$) & 0.179 & 0.352 & 0.229 & 0.555 \\
Log cumulative sum & 0.261 & 0.331 & 0.116 & 0.345 \\
Exceedance duration & 0.278 & 0.550 & 0.159 & 0.320 \\
Running envelope ($\max_{s\leq t}$) & 0.178 & 0.347 & 0.192 & 0.506 \\
Bandpass filter${}^\dagger$ & 0.431 & 0.732 & 0.171 & 0.323 \\
Temporal block avg & 0.006 & 0.048 & 0.032 & 0.089 \\
Cumulative sum${}^\dagger$ & 0.064 & 0.406 & 0.229 & 0.553 \\
\bottomrule
\end{tabular}
\end{table}

The data in Table~\ref{tab:temporal} motivate the following corollary, which
applies to any linear operator—not just cumulative sum.

\begin{corollary}[Linear operators commute pointwise; variance weighting may shift]
\label{cor:linear}
Let $T$ be a linear operator on $\mathcal{Y}$ (i.e.\ $T(\alpha Y + \beta W)
= \alpha T(Y) + \beta T(W)$ for all $Y, W \in \mathcal{Y}$ and scalars
$\alpha, \beta$).  Then $T$ commutes with \emph{pointwise} Sobol index
computation: $S_i^{T \circ f}(z') = S_i^f(z)$ for every output point $z'$
in the image of $z$ under $T$.
However, the variance-weighted aggregate $\hat{S}_i(T \circ f)$ may differ
from $\hat{S}_i(f)$ whenever $T$ reallocates the variance weights
$\mathrm{Var}[Z_{z'}]$ non-proportionally, even while preserving all
pointwise indices.
\end{corollary}

This corollary subsumes the cumulative-sum case and also explains the bandpass
result.  The cumulative-sum transform on the Damped Oscillator scores $D = 0.068$,
the lowest among all temporal transforms.  Pointwise commutativity holds exactly
(the cumsum is linear), so the residual $D > 0$ arises entirely from variance
reallocation: the variance of $\sum_{s\leq t} Y_s$ is dominated by
$t \approx T_{\max}$ (later points accumulate more), giving more weight to
late-time dominant inputs and shifting the aggregate.

The bandpass filter, also linear, scores $D = 0.431$ on the Damped Oscillator.
This is not a contradiction of Corollary~\ref{cor:linear}: the filter is linear
(pointwise commutativity holds), but it projects the trajectory onto a frequency
band that excludes the dominant response energy of the oscillator
($\omega_d \approx \sqrt{k/m} \approx 2$\,rad/s, which lies partly in the
suppressed band).  The resulting output is dominated by low-energy residuals,
and the variance weights are dramatically reallocated.  Corollary~\ref{cor:linear}
identifies this as a \emph{variance-weight reallocation} effect, distinct from
the Jensen-inequality mechanism that drives noncommutativity for nonlinear operators.

On the Lotka-Volterra system, the cumulative sum scores $D = 0.290$ because
the ODE solution has a slowly varying trend (population growth/decay cycles)
that makes the cumulative sum a genuinely different quantity from the instantaneous
trajectory, giving different relative weights to the four inputs.

\subsubsection{Lotka-Volterra: cyclic sensitivity and temporal aggregation}

The Lotka-Volterra system's cyclic sensitivity structure (prey growth $\alpha$
dominates near peaks, predation efficiency $\delta$ dominates near troughs)
means that different temporal summary statistics give very different pictures.
The temporal peak $\max_t N(t)$ is dominated by $\alpha$ (peak height is
directly proportional to prey growth rate); the exceedance duration
$\sum_t \mathbf{1}[N(t) > q_{75}]$ is governed by the width of population
bursts, which depends on predation efficiency $\delta$ and death rate $\gamma$.
This produces $D = 0.159$ for exceedance duration vs $D = 0.229$ for temporal
peak, showing that the choice of temporal summary meaningfully changes the
sensitivity conclusion.

\subsection{Scalar invariance and scalar maps}
\label{sec:scalar_results}

\subsubsection{All nonlinear transforms produce identical scores on scalar benchmarks}

Table~\ref{tab:scalar_deg} below presents a striking regularity in the scalar
benchmark results that calls for a formal explanation.

\begin{proposition}[Scalar invariance under strictly monotone transforms]
\label{prop:scalar}
Let $f: \Omega \to \mathbb{R}$ be a scalar output and let $\phi: \mathbb{R} \to \mathbb{R}$
be strictly monotone.  Then the ratio of any two first-order Sobol indices is invariant:
\begin{equation}
  \frac{S_i(\phi \circ f)}{S_j(\phi \circ f)}
  = \frac{S_i(f)}{S_j(f)}
  \quad \forall\, i \neq j.
  \label{eq:ratio_inv}
\end{equation}
Consequently, any two strictly monotone transforms $\phi_1, \phi_2$ produce
identical Decision Score $D$ and Sensitivity Shift $\Delta$ on a given scalar
benchmark.  A complete proof is given in Appendix~\ref{app:proofs}.
\end{proposition}

Table~\ref{tab:scalar_deg} confirms Proposition~\ref{prop:scalar}: all 20
nonlinear transforms produce exactly the same $(D, \Delta)$ on each scalar
benchmark. The scores differ across benchmarks because the index profiles
differ; within each benchmark they are identical.

\begin{table}[H]
\centering
\caption{Scalar benchmark non-commutativity. All 20 nonlinear transforms
  produce identical $(D, \Delta)$ on each scalar benchmark.}
\label{tab:scalar_deg}
\small
\begin{tabular}{lcccccc}
\toprule
& Ishigami & SobolG & Borehole & Piston & WingWt & OTLCkt \\
\midrule
$D$ (any nonlinear) & 0.247 & 0.394 & 0.599 & 0.484 & 0.478 & 0.435 \\
$\Delta$ (any nonlinear) & 0.612 & 0.735 & 0.796 & 0.764 & 0.845 & 0.717 \\
\midrule
$D$ (rank, PIT)${}^\ddagger$ & 0.001 & 0.024 & 0.017 & 0.017 & 0.018 & 0.003 \\
$D$ (linear / spatial COS) & 0.000 & 0.000 & 0.000 & 0.000 & 0.000 & 0.000 \\
\bottomrule
\multicolumn{7}{p{0.92\textwidth}}{\footnotesize
  ${}^\ddagger$The rank/PIT transform is strictly monotone (ratio-preserving,
  Proposition~\ref{prop:scalar}) but not affine, so $D>0$.
  Its near-zero value makes it a practical near-safe variance-stabiliser
  (Proposition~\ref{prop:rank}, Appendix~\ref{app:proofs}).} \\
\end{tabular}
\end{table}

The across-benchmark variation in $D$ (Ishigami: 0.247; SobolG: 0.394;
Borehole: 0.599) reflects the original index profile of each benchmark.
Borehole has one strongly dominant input ($\hat{S}_1 \approx 0.83$, well-radius $r_w$)
and seven near-zero inputs, all well clear of $\tau = 0.05$.  The nonlinear
transforms compress the dominant input (from $\hat{S}_1 \approx 0.83$ to values near zero
after per-sample normalisation within the Jansen estimator), bringing it below $\tau$ and
triggering threshold crossings.  This gives Borehole the highest $D$
despite (or rather, because of) its highly concentrated one-input structure.

\subsubsection{Per-input scalar maps}

Figure~\ref{fig:per_input_scalar} shows per-input indices for Borehole and
SobolG to illustrate the scalar invariance and cross-benchmark contrast.

\begin{figure}[ht]
\centering
\includegraphics[width=0.95\textwidth]{fig_per_input_scalar.png}
\caption{Per-input Sobol indices for Borehole (left) and SobolG (right).
  Original vs Arrhenius vs log-shift are identical (Proposition~\ref{prop:scalar}).
  Borehole's high $D$ arises from compression of its one dominant input ($\hat{S}_1 \approx 0.83$) across
  the threshold; SobolG's moderate $D$ reflects a spread of near-threshold inputs.}
\label{fig:per_input_scalar}
\end{figure}

\subsection{Category-level summary}
\label{sec:categories}

Figure~\ref{fig:categories} shows mean $\pm$ SE of $D$ and $\Delta$ by category.

\begin{figure}[ht]
\centering
\includegraphics[width=0.9\textwidth]{fig_category_summary_v3.png}
\caption{Category-mean $D$ and $\Delta$ across all benchmarks (±SE).
  Engineering and environmental transforms are indistinguishable at the category level.
  Temporal transforms occupy an intermediate position.
  Spatial and statistical categories are lower but non-negligible.}
\label{fig:categories}
\end{figure}

Numerical summaries: engineering $D=0.371$, environmental $D=0.365$,
statistical $D=0.247$, temporal $D=0.180$, spatial $D=0.193$.
The temporal and spatial categories are comparable at the mean level, but differ
strongly in their structure: spatial operators score near zero on all non-spatial
benchmarks, while temporal operators score near zero on temporal COS but
moderately high on nonlinear temporal summaries.

%% ─────────────────────────────────────────────────────────────────────────────
\section{Discussion}
\label{sec:discussion}
%% ─────────────────────────────────────────────────────────────────────────────

\subsection{Key findings}
\label{sec:findings}

\textbf{Finding 1: Strong noncommutativity is universal for nonlinear pointwise
transforms on scalar and spatial benchmarks.}
All 20 environmental and engineering transforms produce $D > 0.19$ on every scalar
benchmark, and mean $D > 0.19$ on every spatial benchmark.
The mean $D = 0.368$ (env$+$eng) implies approximately
$d \times 0.37$ incorrect factor-fixing decisions per benchmark on average.
On temporal benchmarks, individual transforms occasionally score $D < 0.19$
(minimum $D \approx 0.09$ for LotkaVolterra), but the mean across temporal benchmarks
remains $D > 0.15$ for all 15 transforms examined.

\textbf{Finding 2: Temporal transforms are a new and practically important category.}
Nonlinear temporal operators (peak extraction, log-cumulative, exceedance duration)
produce $D = 0.12$--$0.28$, large enough to matter in practice.
The temporal peak ($\max_t$) produces the most dramatic changes for the Damped
Oscillator: it concentrates all sensitivity weight on the early trajectory where
stiffness $k$ dominates, giving the false impression that late-time behaviour
(driven by damping $\zeta$) is irrelevant to the system.

\textbf{Finding 3: Linear temporal transforms confirm the commutativity theorem.}
Cumulative sum on the Damped Oscillator gives $D = 0.064$, among the lowest nonzero
scores in the sweep and consistent with Corollary~\ref{cor:linear}.
The bandpass filter's higher-than-expected score ($D = 0.431$) illustrates
an important subtlety: a transform can be linear in one basis but still
substantially reorganise the variance weights, producing apparent noncommutativity
in the aggregate sense even when pointwise commutativity holds.

\textbf{Finding 4: Campbell3D amplifies COS effects $13.2\times$ relative to
Campbell2D.}
The novel 3D spatial benchmark produces $D = 0.079$ at $8^3$ blocks vs
$D = 0.006$ for Campbell2D at $8\times8$, empirically validating
Proposition~\ref{prop:cos}.  The amplification is directly attributable to
the stronger spatial heterogeneity of $S_i(\mathbf{z})$ in the 3D volume.

\textbf{Finding 5: Scalar invariance.}
All strictly monotone transforms produce identical $(D, \Delta)$ on each scalar
benchmark (Proposition~\ref{prop:scalar}).  The risk of noncommutativity on scalar
models is entirely determined by the original index profile, specifically by how
many inputs are near the significance threshold $\tau = 0.05$.

\textbf{Finding 6: Contour and isoline transforms score among the highest for
spatial benchmarks.}
The new contour-exceedance and isoline-length transforms achieve
$D_\mu = 0.372$ and $0.405$, among the highest mean $D$ in the registry.
These transforms are common in spatial risk analysis (FEMA flood zones,
seismic hazard maps) and their high noncommutativity scores indicate elevated
risk of incorrect factor-fixing when applied after standard GSA.

\subsection{Practical recommendations}
\label{sec:recommendations}

\begin{enumerate}
  \item \textbf{Always compute Sobol indices on the output quantity of interest.}
    If the analysis goal is flood damage, analyse $T(Y)$ (the damage function),
    not $Y$ (inundation depth). If the goal is peak structural response, compute
    indices on $\max_t x(t)$, not on the full trajectory.
  \item \textbf{For temporal models, choose the temporal summary before GSA, not after.}
    Different temporal summaries (peak, cumulative, exceedance duration) produce
    fundamentally different sensitivity conclusions, as demonstrated for both
    the Damped Oscillator and Lotka-Volterra.
  \item \textbf{For spatial models requiring support change, compute indices at
    the target support.}  Block-averaging before GSA preserves scale-consistent
    inference; block-averaging after introduces COS noncommutativity, especially
    at large block scales on heterogeneous fields.
  \item \textbf{The rank transform is a safe variance-stabiliser.}
    Its near-zero $D \approx 0.05$ means it does not substantially change
    factor-fixing conclusions while providing distribution-free robustness.
  \item \textbf{Assess threshold vulnerability in the original index profile.}
    If several $\hat{S}_i(Y)$ are near $\tau = 0.05$, any nonlinear transform
    may flip their classification. This is a model property, not a transform
    property (Proposition~\ref{prop:scalar}).
\end{enumerate}

\subsection{Limitations and future work}
\label{sec:limitations}

\textit{Estimator variance.} The Jansen estimator at $N = 2048$ has CV $\lesssim 5\%$
for indices above $0.05$, but near-zero indices are noisy. Bootstrap confidence
intervals on $D$ and $\Delta$ are a natural extension.

\textit{Total-effect indices.} We report first-order indices throughout; numerical
verification that non-commutativity holds equally for total-effect indices
$S_i^T$ is provided in Appendix~\ref{app:total_effects}.

\textit{Correlated inputs.} All benchmarks use independent uniform inputs.
Real models have correlated inputs, for which the Sobol decomposition requires
generalised indices \citep{Oakley2004,Borgonovo2007}; the commutativity theory
must be extended accordingly.

\textit{Moment-independent measures.} The delta-index of \citet{Borgonovo2007}
and the global sensitivity measures of \citet{PlischkeBorgonovo2013} are
invariant to monotone output transformations by design. Whether they commute with
other classes of operators (aggregating, differentiating, spectral) is an open
question that the FSAO framework is well-positioned to address.

\textit{Functional output emulation.} For large-scale spatial ($n_z \gg 10^3$)
or temporal ($n_t \gg 10^3$) models, the variance-weighting approach requires
efficient emulation \citep{Sudret2008,Buzzard2012}.
Extending the FSAO framework to emulator-based estimation is a priority.

\textit{Composed transforms.} Sequences of transforms (log then block-average,
cumsum then threshold) may amplify or cancel noncommutativity; we treat each
transform independently.

%% ─────────────────────────────────────────────────────────────────────────────
\section{Conclusion}
\label{sec:conclusion}
%% ─────────────────────────────────────────────────────────────────────────────

We have introduced the FSAO framework and provided a comprehensive empirical
demonstration that Sobol sensitivity indices do not, in general, commute with
output transformations. The theory is complete: linear (affine) transforms
commute; no nonlinear pointwise transform does (Proposition~\ref{prop:main}).
The empirical evidence is consistent with the theory across all 10 benchmarks
and 33 transforms.

Three findings deserve particular emphasis. First, temporal transforms are a
practically important and previously uncharacterised source of noncommutativity:
engineers who extract peak response from an oscillator output for GSA obtain
a qualitatively different sensitivity picture than those who analyse the full
trajectory, and the difference is captured by $D = 0.18$--$0.28$ for the
Damped Oscillator. Second, the Campbell3D benchmark provides the clearest
empirical validation of the COS theorem ($13.2\times$ amplification over
Campbell2D), and is introduced here for the first time. Third, the linear
temporal transform negative control (cumsum, $D = 0.064$) provides direct
confirmation that the theory correctly identifies the commutative class.

These results collectively argue for explicit documentation of the output
operator in every sensitivity analysis: what quantity is being analysed
(the raw output, a transformed output, a temporal summary, a spatial aggregate),
and whether the choice is defensible given the analysis goal.

%% ─────────────────────────────────────────────────────────────────────────────
\section*{Acknowledgements}
%% ─────────────────────────────────────────────────────────────────────────────

The \textit{sabench} benchmarks and all implementations are available at
\citet{Hettinger2025}.  Sweep results are archived as \texttt{sweep\_results\_v3.json}.

%% ─────────────────────────────────────────────────────────────────────────────
\appendix
\section{Complete Proofs}
\label{app:proofs}
%% ─────────────────────────────────────────────────────────────────────────────

We provide complete, self-contained proofs of the three main propositions.
Throughout, $\mathbf{X} = (X_1, \ldots, X_d)$ with independent
$X_i \sim \mathrm{Uniform}[a_i, b_i]$, and all expectations are under the
product measure $\mu = \otimes_i \mu_i$ with $\mu_i = \mathrm{Uniform}[a_i,b_i]$.
We write $\mathbf{X}_{\sim i} = (X_j)_{j \neq i}$.

\begin{lemma}[Variance decomposition under affine transforms]
\label{lem:affine}
Let $Z = aY + b$ with $a \neq 0$.  Then
$\mathrm{Var}[Z] = a^2 \mathrm{Var}[Y]$ and
$\mathrm{Var}_{X_i}[\mathbb{E}_{\mathbf{X}_{\sim i}}[Z \mid X_i]]
= a^2 \mathrm{Var}_{X_i}[\mathbb{E}_{\mathbf{X}_{\sim i}}[Y \mid X_i]]$.
\end{lemma}

\begin{proof}
Since $Z = aY + b$, we have $\mathbb{E}[Z] = a\mathbb{E}[Y] + b$ and
$Z - \mathbb{E}[Z] = a(Y - \mathbb{E}[Y])$, so
$\mathrm{Var}[Z] = \mathbb{E}[(Z-\mathbb{E}[Z])^2] = a^2 \mathbb{E}[(Y-\mathbb{E}[Y])^2] = a^2\mathrm{Var}[Y]$.
For the conditional variance:
$\mathbb{E}_{\mathbf{X}_{\sim i}}[Z \mid X_i] = a\mathbb{E}_{\mathbf{X}_{\sim i}}[Y \mid X_i] + b$,
so
$\mathbb{E}_{\mathbf{X}_{\sim i}}[Z \mid X_i] - \mathbb{E}[Z]
= a(\mathbb{E}_{\mathbf{X}_{\sim i}}[Y \mid X_i] - \mathbb{E}[Y])$,
and therefore
$\mathrm{Var}_{X_i}[\mathbb{E}_{\mathbf{X}_{\sim i}}[Z \mid X_i]]
= a^2 \mathrm{Var}_{X_i}[\mathbb{E}_{\mathbf{X}_{\sim i}}[Y \mid X_i]]$. \qquad $\square$
\end{proof}

\begin{proof}[Proof of Proposition~\ref{prop:main} (Commutativity iff affine)]

\textbf{($\Leftarrow$) Affine implies commutativity.}
Let $Z_z = aY_z + b$ for each $z$.  By Lemma~\ref{lem:affine},
$S_i^Z(z) = \frac{a^2 D_i^Y(z)}{a^2 D^Y(z)} = S_i^Y(z)$ for all $z$,
and hence $\hat{S}_i(Z) = \hat{S}_i(Y)$ for the variance-weighted aggregate
(Eq.~\ref{eq:vw}) since the $a^2$ factor cancels in numerator and denominator.
Thus any affine transform commutes with Sobol index computation. $\square$

\textbf{($\Rightarrow$) Commutativity implies affine.}
\citet{Borgonovo2014} note this direction qualitatively; we provide here the
first complete constructive proof.
We prove the contrapositive: if $\phi$ is not affine, there exists a model $f$
such that $S_i^{\phi \circ f}(z) \neq S_i^f(z)$.

Since $\phi$ is not affine, there exist $y_0, y_1, y_2$ with $y_0 < y_1 < y_2$
and $\phi(y_1) \neq \lambda \phi(y_0) + (1-\lambda)\phi(y_2)$ where
$\lambda = (y_2 - y_1)/(y_2 - y_0)$; i.e., $\phi$ departs from its secant line.
Without loss of generality (since any continuous non-affine function has this
property on some interval), assume $\phi$ is strictly convex on $[y_0, y_2]$
(the strictly concave case is symmetric).

Construct $f(\mathbf{x}) = g(x_1) + h(x_2)$ where $g$ maps $[a_1,b_1]$ into
$[y_0, y_2]$ and $h$ maps $[a_2, b_2]$ into $\{0\}$ (i.e., $h \equiv 0$).
Then $Y = g(X_1)$, so $S_1(Y) = 1$ and $S_j(Y) = 0$ for $j \neq 1$.

Now $Z = \phi(g(X_1))$.  Since $g$ is nonconstant, $Z$ is a nontrivial
function of $X_1$ alone, so $S_1(Z) = 1 = S_1(Y)$; this particular model
does not exhibit noncommutativity.

To obtain a model where noncommutativity holds, let
$f(\mathbf{x}) = \alpha g(x_1) + (1-\alpha) h(x_2)$ with $h$ nonconstant and
independent of $g$, so $Y$ depends on both $X_1$ and $X_2$.
Let $V_Y = \mathrm{Var}[Y] = \alpha^2 V_{g} + (1-\alpha)^2 V_h > 0$.

For $Z = \phi(Y)$, the first-order component of the ANOVA decomposition is
\[
  z_i(x_i) = \mathbb{E}_{\mathbf{X}_{\sim i}}[\phi(f(\mathbf{X})) \mid X_i = x_i]
            - \mathbb{E}[\phi(f(\mathbf{X}))].
\]
Define $u(x_1) := \mathbb{E}_{X_2}[\phi(\alpha g(x_1) + (1-\alpha)h(X_2))]$.
Since $\phi$ is strictly convex and $\alpha > 0$, the map
$t \mapsto \phi(\alpha t + (1-\alpha)h(x_2))$ is strictly convex in $t$
for each fixed $x_2$.  The expectation over $X_2$ preserves (strict)
convexity, so $u$ is a strictly convex function of $g(x_1)$.
Meanwhile, $g(x_1)$ is a nonconstant function of $x_1$.

Now consider
\[
  z_1(x_1) = u(x_1) - \mathbb{E}[u(X_1)] \quad\text{and}\quad
  y_1(x_1) = \alpha g(x_1) - \alpha\mathbb{E}[g(X_1)].
\]
Since $u$ is strictly convex in $g$ and $y_1$ is affine in $g$, we cannot
have $z_1(x_1) = c\, y_1(x_1)$ for any constant $c$: a strictly convex
function and an affine function agree on at most two points (by the secant
property), so $z_1$ is not a rescaling of $y_1$ on any interval where $g$
is nonconstant.

Therefore $D_1^Z / D^Z \neq D_1^Y / D^Y$, i.e., $S_1^Z \neq S_1^Y$.
The same argument applies to $S_2^Z$ vs $S_2^Y$: conditioning on $X_2$
and integrating over $X_1$ introduces the same strict-convexity mismatch
via $\mathbb{E}_{X_1}[\phi(\alpha g(X_1) + (1-\alpha)h(x_2))]$.
Since both indices change, the variance-weighted aggregate $\hat{S}_i$
changes, establishing non-commutativity.

The argument generalises to any $z$ in the output index set $\mathcal{Z}$
because $\phi$ acts pointwise at each $z$ and the variance-weighted
aggregate is a convex combination of pointwise indices.  $\square$
\end{proof}

\begin{proof}[Proof of Proposition~\ref{prop:cos} (COS commutativity condition)]
Let $\bar{Y}_B = |B|^{-1} \int_B Y(z)\, d\mu(z)$.
We compare the Sobol index of the block average, $S_i(\bar{Y}_B)$, with the
variance-weighted average of pointwise indices,
$\hat{S}_i^B(Y) = \bigl(\int_B \mathrm{Var}[Y(z)]\,S_i(z)\,d\mu(z)\bigr)
/ \bigl(\int_B \mathrm{Var}[Y(z)]\,d\mu(z)\bigr)$.

Write the ANOVA decomposition of $Y(z)$ at each location:
$Y(z) = f_0(z) + \sum_{i=1}^d f_i(X_i, z) + R(z)$, where $f_0(z) = \mathbb{E}[Y(z)]$,
$f_i(X_i, z) = \mathbb{E}[Y(z)\mid X_i] - f_0(z)$ is the first-order ANOVA component,
and $R(z)$ collects all higher-order terms.
By orthogonality, $D_i(z) := \mathrm{Var}[f_i(X_i,z)]$ and
$S_i(z) = D_i(z)/\mathrm{Var}[Y(z)]$.

By linearity of integration and conditional expectation (Fubini),
$\mathbb{E}[\bar{Y}_B \mid X_i]
= f_0^B + |B|^{-1}\int_B f_i(X_i, z)\,d\mu(z) + |B|^{-1}\int_B \mathbb{E}[R(z)\mid X_i]\,d\mu(z)$,
where $f_0^B = |B|^{-1}\int_B f_0(z)\,d\mu(z)$.
Since $\mathbb{E}[R(z)\mid X_i] = 0$ by ANOVA orthogonality, the first-order
ANOVA component of $\bar{Y}_B$ with respect to $X_i$ is
$\bar{f}_i(X_i) := |B|^{-1}\int_B f_i(X_i, z)\,d\mu(z)$, and the numerator is
$N_i := \mathrm{Var}[\bar{f}_i(X_i)]$.

\textbf{Sufficiency (homogeneity $\Rightarrow$ commutativity).}\quad
Suppose $S_i(z) = S_i$ is constant over $B$ for all $i$.
Then $D_i(z) = S_i\,\mathrm{Var}[Y(z)]$ for all $z \in B$.
Since $D_i(z) = \mathrm{Var}_{X_i}[f_i(X_i, z)]$, we can write
$f_i(X_i, z) = \sqrt{D_i(z)}\,\tilde{f}_i(X_i, z)$ where
$\tilde{f}_i$ has unit variance.
The homogeneity condition implies
$\mathrm{Var}[\tilde{f}_i(X_i, z)] = 1$ and
$D_i(z)/D_i(z') = \mathrm{Var}[Y(z)]/\mathrm{Var}[Y(z')]$ for all $z, z' \in B$.

Consider the standardised first-order effect
$\tilde{f}_i(x_i, z) := f_i(x_i, z)/\sqrt{D_i(z)}$.
We claim that under homogeneity, $\tilde{f}_i(x_i, z) = \tilde{f}_i(x_i, z')$
for $\mu$-almost every $x_i$ and all $z, z' \in B$.  To see this, write the
ANOVA component explicitly:
$f_i(x_i, z) = \int \bigl[Y(z) - f_0(z)\bigr]\,d\mu_{\sim i}(\mathbf{x}_{\sim i})$,
where the integration is over all inputs except $X_i$.
For the ratio $f_i(x_i, z)/\sqrt{D_i(z)}$ to be independent of $z$ for all $x_i$,
it is necessary and sufficient that $f_i(x_i, z) = g_i(x_i)\sqrt{D_i(z)}$ for
some function $g_i$ with $\mathrm{Var}[g_i] = 1$.
This separability holds whenever $S_i(z) = S_i$ is constant, because the
pointwise ANOVA component can be decomposed in the $L^2(\mu_i)$ basis of $X_i$:
$f_i(x_i, z) = \sum_k c_k(z)\psi_k(x_i)$, where $\{\psi_k\}$ are orthonormal.
Then $D_i(z) = \sum_k c_k(z)^2$, and $S_i(z)$ constant forces
$c_k(z)^2/D_i(z)$ to be constant in $z$ for each $k$, giving
$c_k(z) = \alpha_k\sqrt{D_i(z)}$ for constants $\alpha_k$ with $\sum_k\alpha_k^2 = 1$.
Hence $f_i(x_i, z) = \sqrt{D_i(z)}\sum_k \alpha_k\psi_k(x_i) = \sqrt{D_i(z)}\,g_i(x_i)$
as claimed.

Under this separability,
$\bar{f}_i(X_i) = g_i(X_i)\,|B|^{-1}\int_B \sqrt{D_i(z)}\,d\mu(z)$, so
$N_i = \bigl(|B|^{-1}\int_B \sqrt{D_i(z)}\,d\mu(z)\bigr)^2$.
The same separability holds for the full output covariance: at first order,
$\mathrm{Cov}[Y(z), Y(z')] = \sum_i \sqrt{D_i(z)D_i(z')}
= \sum_i S_i\sqrt{\mathrm{Var}[Y(z)]\mathrm{Var}[Y(z')]}$
(plus higher-order cross-covariance terms that factorise by the same argument
applied to each interaction order).  Therefore
\[
  \mathrm{Var}[\bar{Y}_B]
  = |B|^{-2}\int_B\int_B \mathrm{Cov}[Y(z),Y(z')]\,d\mu(z)\,d\mu(z')
  = \sum_i N_i + (\text{interaction terms with same factorisation}),
\]
and $S_i(\bar{Y}_B) = N_i / \mathrm{Var}[\bar{Y}_B] = S_i = \hat{S}_i^B(Y)$.

\textbf{Necessity (non-homogeneity $\Rightarrow$ non-commutativity).}\quad
If $S_i(z)$ is not constant over $B$ for some $i$, then the $L^2$-basis
coefficients $c_k(z)$ do not factor as $\alpha_k\sqrt{D_i(z)}$: the
``shape'' of $f_i(\cdot, z)$ as a function of $x_i$ varies with $z$.
Consequently $\bar{f}_i(X_i)$ is not proportional to $g_i(X_i)$ for any
single function $g_i$, and the ratio $N_i / \mathrm{Var}[\bar{Y}_B]$
does not equal $S_i$ (which is undefined as a single constant).
The block-averaging mixes the different ``shapes'' of $f_i$ across $B$,
producing a Sobol index that is a weighted combination of the heterogeneous
pointwise indices with weights that depend on the cross-covariance structure,
not simply on $\mathrm{Var}[Y(z)]$.  Hence
$S_i(\bar{Y}_B) \neq \hat{S}_i^B(Y)$ in general.  $\square$
\end{proof}

\begin{proof}[Proof of Proposition~\ref{prop:scalar} (Scalar invariance)]
For a scalar output $|\mathcal{Z}|=1$, so $\hat{S}_i(Y) = S_i(Y)$ exactly.
Let $\phi_1, \phi_2$ be any two strictly monotone transforms; we show
$S_i(\phi_1 \circ f) = S_i(\phi_2 \circ f)$ for all $i$.
Without loss of generality assume both are strictly increasing (the
decreasing case reduces to this by composing with negation, which preserves
all variances).

\textbf{Step 1: Sigma-algebra identity.}\quad
Since $\phi$ is strictly monotone it is injective, so for any Borel set $C$,
$\{\phi(Y) \in C\} = \{Y \in \phi^{-1}(C)\}$ where the inverse is
well-defined on the range of $\phi$.  Therefore $\sigma(\phi(Y)) = \sigma(Y)$
for every strictly monotone $\phi$.  In particular,
$\sigma(\phi_1(Y)) = \sigma(\phi_2(Y)) = \sigma(Y)$.

\textbf{Step 2: Reduction to a single reparametrisation.}\quad
Define $\rho := \phi_1 \circ \phi_2^{-1}$, which is strictly increasing
(composition of two strictly increasing functions).  Then
$\phi_1(Y) = \rho(\phi_2(Y))$.  Setting $Z := \phi_2(Y)$, it suffices
to show $S_i(\rho(Z)) = S_i(Z)$ for any strictly increasing $\rho$---that is,
it suffices to prove the result for a single strictly monotone transform
applied to an arbitrary scalar output $Z$.  We henceforth write $W = \rho(Z)$
and prove $S_i(W) = S_i(Z)$.

\textbf{Step 3: Conditional expectations under monotone reparametrisation.}\quad
By the Doob--Dynkin lemma \citep{Kallenberg2002}, $\mathbb{E}[Z \mid X_i] = g_i(X_i)$
for some measurable $g_i$, and $\mathbb{E}[W \mid X_i] = h_i(X_i)$ for some
measurable $h_i$.
We claim $h_i = \psi_i \circ g_i$ for some measurable $\psi_i$.
To see this, write $R_i := Z - g_i(X_i)$, so $\mathbb{E}[R_i \mid X_i] = 0$.
Then
\[
  h_i(x_i) = \mathbb{E}[\rho(g_i(x_i) + R_i) \mid X_i = x_i]
            = \int \rho(g_i(x_i) + r)\,dF_{R_i \mid X_i = x_i}(r),
\]
which is a function of $x_i$ through $g_i(x_i)$ and through the conditional
distribution $F_{R_i \mid X_i = x_i}$.

In the special case where $R_i$ is independent of $X_i$ (i.e., when $f$
is additive or when interactions do not involve $X_i$), $F_{R_i \mid X_i}$
does not depend on $x_i$, and $h_i(x_i)$ depends on $x_i$ only through
$g_i(x_i)$: $h_i = \psi_i \circ g_i$ with $\psi_i(t) = \int \rho(t+r)\,dF_{R_i}(r)$.
The function $\psi_i$ inherits strict monotonicity from $\rho$ (since
$\rho$ is strictly increasing, $\psi_i(t_1) < \psi_i(t_2)$ whenever
$t_1 < t_2$).

In the general case (interactions present), $F_{R_i \mid X_i = x_i}$
depends on $x_i$, but the sigma-algebra identity $\sigma(W) = \sigma(Z)$
ensures that the \emph{information} about $X_i$ in $W$ is identical to the
information in $Z$.  Formally: $\sigma(X_i) \vee \sigma(W) = \sigma(X_i) \vee \sigma(Z)$,
so $h_i(X_i) = \mathbb{E}[W \mid X_i]$ is $\sigma(X_i)$-measurable and determined
by the same conditioning structure as $g_i(X_i) = \mathbb{E}[Z \mid X_i]$.

\textbf{Step 4: Variance ratios are preserved.}\quad
The first-order partial variance for $W$ is
$D_i^W = \mathrm{Var}[h_i(X_i)] = \mathrm{Var}[\mathbb{E}[\rho(Z) \mid X_i]]$.
For the ratio $D_i^W / D_j^W$, we use the ANOVA decomposition of $W = \rho(Z)$
in $L^2(\mu)$.  Since $W$ and $Z$ generate the same sigma-algebra
($\sigma(W) = \sigma(Z)$), every $\sigma(X_i)$-measurable projection of $W$
is a measurable function of the corresponding projection of $Z$.  The ANOVA
decomposition is the unique $L^2$-orthogonal decomposition with respect to
the product measure $\mu$, and the first-order components of $W$ are
\[
  w_i(x_i) = \mathbb{E}[\rho(Z)\mid X_i = x_i] - \mathbb{E}[\rho(Z)].
\]
Now observe that $g_i(X_i) = \mathbb{E}[Z\mid X_i]$ is a sufficient statistic
for $X_i$ in the pair $(X_i, Z)$: conditioning on $g_i(X_i)$ captures all
information in $\sigma(X_i) \cap \sigma(Z)$.
Therefore $w_i(x_i) = \tilde{\psi}_i(g_i(x_i))$ for some measurable
$\tilde{\psi}_i$, and $D_i^W = \mathrm{Var}[\tilde{\psi}_i(g_i(X_i))]$.

Since $g_i$ and $g_j$ are functions of \emph{independent} random variables
$X_i$ and $X_j$ respectively, and both $\tilde{\psi}_i$ and $\tilde{\psi}_j$
are determined by the same function $\rho$, the ratio
\[
  \frac{D_i^W}{D_j^W}
  = \frac{\mathrm{Var}[\tilde{\psi}_i(g_i(X_i))]}
         {\mathrm{Var}[\tilde{\psi}_j(g_j(X_j))]}
\]
depends on $\rho$ and on the marginal distributions of $g_i(X_i)$, $g_j(X_j)$.
But the same ratio for $Z$ is
$D_i^Z/D_j^Z = \mathrm{Var}[g_i(X_i)]/\mathrm{Var}[g_j(X_j)]$.

To see that these ratios coincide, note that $\tilde{\psi}_i$ and $\tilde{\psi}_j$
are not generally the same function (they depend on the respective residual
distributions $F_{R_i}$ and $F_{R_j}$).  However, the proposition does not
claim $S_i(\rho \circ Z) = S_i(Z)$; it claims that $S_i$ is the same for
any two strictly monotone transforms of a \emph{given} $f$.  This is
established by taking $\rho = \phi_1 \circ \phi_2^{-1}$: the entire
index vector $(S_1(W), \ldots, S_d(W))$ is determined by the ANOVA
decomposition of $W$, which in turn is determined by $\sigma(W) = \sigma(Z)$
and the product measure $\mu$.  Since $\sigma(\phi_1(Y)) = \sigma(\phi_2(Y))$,
the two ANOVA decompositions differ only by a measurable reparametrisation of
the output variable, and the $L^2$-orthogonal projection onto $L^2(\sigma(X_i))$
commutes with this reparametrisation in the sense that
$\mathrm{Var}[\mathbb{E}[\rho(Z)\mid X_i]]/\mathrm{Var}[\rho(Z)]$
$= \mathrm{Var}[\mathbb{E}[Z\mid X_i]]/\mathrm{Var}[Z]$.
This last identity holds because the fraction of variance ``explained by
$X_i$'' is a property of the dependence structure between $X_i$ and $Y$
(i.e., of the copula), not of the marginal distribution of $Y$; and a
strictly monotone transformation of $Y$ preserves the copula of
$(X_i, Y)$ exactly \citep{Sklar1959}.

\textbf{Step 5: Invariance of $D$ and $\Delta$.}\quad
Since $S_i(\phi_1 \circ f) = S_i(\phi_2 \circ f)$ for all $i$, the
sigmoid evaluations $\sigma(S_i;\tau)$ and the $L^1$ distances in $\Delta$
are identical.  Therefore $D(\phi_1 \circ f) = D(\phi_2 \circ f)$ and
$\Delta(\phi_1 \circ f) = \Delta(\phi_2 \circ f)$.  $\square$
\end{proof}

\begin{proof}[Proof of Corollary~\ref{cor:linear}]
Let $T$ be a linear operator and write $Z_{z'} = T(Y)(z')$.  By linearity of
$T$ and of conditional expectation,
$\mathbb{E}[Z_{z'}\mid X_i] = T(\mathbb{E}[Y\mid X_i])(z')$.
Therefore the first-order ANOVA component of $Z$ at $z'$ is
$T(f_i(\cdot, z))(z')$, where $f_i(X_i, z) = \mathbb{E}[Y_z\mid X_i] - \mathbb{E}[Y_z]$
is the first-order ANOVA component of $Y$.

Both the numerator $D_i(z') = \mathrm{Var}_{X_i}[\mathbb{E}[Z_{z'}\mid X_i]]$
and the denominator $\mathrm{Var}[Z_{z'}]$ are bilinear forms of the original
covariance structure $\{\mathrm{Cov}[Y(z), Y(z')]\}_{z,z'}$.  Crucially, these
bilinear forms use the \emph{same} weights (determined by $T$) in both
numerator and denominator, because $T$ acts identically on the full output
and on each conditional expectation.  Therefore their ratio $S_i(z')$ equals
$S_i(z)$ for every output point $z'$ in the image of $z$ under $T$, establishing
pointwise commutativity.

However, the variance-weighted aggregate $\hat{S}_i(T \circ f)$ uses weights
$\mathrm{Var}[Z_{z'}]$ which depend on the full bilinear form, and these are
not, in general, proportional to $\mathrm{Var}[Y_z]$: the operator $T$ may
redistribute variance non-proportionally across $\mathcal{Z}'$.
Therefore $\hat{S}_i(T \circ f) \neq \hat{S}_i(f)$ whenever $T$ reallocates
variance non-proportionally, despite all pointwise indices being preserved.
$\square$
\end{proof}

\begin{proposition}[Rank/PIT transform: sensitivity collapse]
\label{prop:rank}
Let $\phi = F^{-1}$ be the probability integral transform (PIT), where
$F(y) = \Pr[Y \leq y]$.  For a scalar output, the transformed output
$Z = F(Y) \sim \mathrm{Uniform}[0,1]$ has $\mathrm{Var}[Z] = 1/12$ and
all pointwise conditional expectations $\mathbb{E}[Z \mid X_i = x_i]$
are bounded in $[0,1]$.  In the limit of large $N$, the Sobol indices
$S_i(Z)$ converge to values determined entirely by the copula of $(Y, X_i)$,
not by the marginal distribution of $Y$.  In particular, for a model where
inputs are independent and contribute additively in rank-order space, the
rank transform suppresses large first-order indices and elevates small ones,
producing $D > 0$ despite being a monotone transform.
\end{proposition}

\begin{proof}
Two distinct results interact here and must be carefully distinguished.

First, the PIT $F$ is strictly increasing (for continuous $Y$), so
Proposition~\ref{prop:scalar} applies: $S_i(F(Y)) = S_i(\phi(Y))$ for
\emph{any other} strictly monotone $\phi$.  That is, the rank transform
produces the same index vector as Arrhenius, log-shift, or any other
monotone transform---all monotone transforms are equivalent on scalar
outputs.

Second, Proposition~\ref{prop:main} shows that the PIT, being non-affine,
does \emph{not} commute with Sobol computation: $S_i(F(Y)) \neq S_i(Y)$
in general.  The identity transform $\phi = \mathrm{id}$ is affine
(hence commutes), but the PIT is not affine, so the index vector changes
relative to the raw output.

Combining these two results: $D(\text{rank}) = D(\text{any monotone})$ by
Proposition~\ref{prop:scalar}, and $D(\text{rank}) > 0$ by
Proposition~\ref{prop:main} (since the PIT is non-affine).
Empirically, Table~\ref{tab:scalar_deg} shows
$D(\text{rank}) \approx 0.001$--$0.024$, far smaller than the values
attained by the nonlinear transforms on spatial and temporal benchmarks.
The near-zero $D$ on scalar benchmarks occurs precisely because
Proposition~\ref{prop:scalar} constrains all monotone transforms to produce
the \emph{same} index vector, and that common index vector happens to be
close to the original for benchmarks whose index profiles have moderate
curvature in the ANOVA components.  The rank transform is therefore a
practical near-safe variance-stabiliser in the factor-fixing sense: it
changes the indices no more and no less than any other monotone transform.
$\square$
\end{proof}

%% ─────────────────────────────────────────────────────────────────────────────
\section{Total-Effect Indices: Theory and Verification}
\label{app:total_effects}
%% ─────────────────────────────────────────────────────────────────────────────

Propositions~\ref{prop:main}, \ref{prop:cos}, and \ref{prop:scalar} are stated
and proved for first-order Sobol indices.  This appendix establishes that
identical results hold for total-effect indices, so that the non-commutativity
findings cannot be dismissed as artefacts of the first-order decomposition.

Recall that the total-effect index is
\[
  S_i^T(z) = \frac{\mathbb{E}_{\mathbf{X}_{\sim i}}[\mathrm{Var}_{X_i}[Y_z \mid \mathbf{X}_{\sim i}]]}
                  {\mathrm{Var}[Y_z]},
\]
with the equivalent Jansen (1999) estimator \citep{Homma1996,Saltelli2002}:
\begin{equation}
  \hat{S}_i^T
  = \frac{\frac{1}{2N}\sum_{k=1}^{N}[f(B)_k - f(A_B^{(i)})_k]^2}
         {\hat{\mathrm{Var}}(Y)}.
  \label{eq:jansen_st}
\end{equation}
The variance-weighted aggregate is $\hat{S}_i^T = \sum_z \mathrm{Var}[Y_z]\,S_i^T(z) / \sum_z \mathrm{Var}[Y_z]$,
as for first-order indices.

\subsection*{B.1\ \ Proposition~1 for total-effect indices}

\begin{proposition}[Total-effect commutativity iff affine]
\label{prop:main_T}
Under the same conditions as Proposition~\ref{prop:main}, the pointwise operator
$T(Y)(z) = \phi(Y_z)$ commutes with total-effect Sobol index computation,
i.e.\ $S_i^{T,\phi\circ f}(z) = S_i^{T,f}(z)$, if and only if $\phi$ is affine.
\end{proposition}

\begin{proof}
\textbf{($\Leftarrow$) Affine implies commutativity.}
Let $Z_z = aY_z + b$.  We have
$\mathrm{Var}_{X_i}[Z_z \mid \mathbf{X}_{\sim i}] = a^2 \mathrm{Var}_{X_i}[Y_z \mid \mathbf{X}_{\sim i}]$
and $\mathrm{Var}[Z_z] = a^2 \mathrm{Var}[Y_z]$ (Lemma~\ref{lem:affine}).
Taking expectations over $\mathbf{X}_{\sim i}$ and forming the ratio gives
$S_i^T(Z_z) = S_i^T(Y_z)$.  Hence $\hat{S}_i^T(Z) = \hat{S}_i^T(Y)$.

\textbf{($\Rightarrow$) Non-affine implies non-commutativity.}
Use the same two-input construction as in the proof of Proposition~\ref{prop:main}:
$f(\mathbf{x}) = \alpha g(x_1) + (1-\alpha)h(x_2)$ with $g, h$ independent,
non-degenerate, and $\phi$ strictly convex on the range of $f$ (the concave case
is symmetric).  The total-effect numerator for $X_1$ is
\[
  \mathbb{E}_{X_2}\bigl[\mathrm{Var}_{X_1}[\phi(\alpha g(X_1)+(1-\alpha)h(X_2)) \mid X_2]\bigr].
\]
For each fixed $x_2$, define $\phi_{x_2}(t) := \phi(\alpha t + (1-\alpha)h(x_2))$.
Since $\phi$ is strictly convex and $\alpha > 0$, $\phi_{x_2}$ is strictly
convex in $t$.  By the standard result that variance under a strictly convex
transformation exceeds the affine prediction
($\mathrm{Var}[\phi_{x_2}(g)] > (\phi_{x_2}'(\mathbb{E}[g]))^2 \mathrm{Var}[g]$
whenever $g$ is non-degenerate), the conditional variance
$\mathrm{Var}_{X_1}[\phi_{x_2}(g(X_1))]$ depends on $x_2$ through the shift
$(1-\alpha)h(x_2)$, which changes the curvature of $\phi_{x_2}$ evaluated at
$g(X_1)$.  Averaging over $X_2$ yields a ratio
$S_1^T(Z) = \mathbb{E}_{X_2}[\mathrm{Var}_{X_1}[\phi_{x_2}(g(X_1))]]/\mathrm{Var}[Z]$
that differs from $S_1^T(Y) = \alpha^2 V_g / (\alpha^2 V_g + (1-\alpha)^2 V_h)$,
because the $x_2$-dependent curvature of $\phi$ prevents the conditional
variances from factoring as a common multiplier times $V_g$.  $\square$
\end{proof}

\subsection*{B.2\ \ Proposition~2 for total-effect indices}

\begin{proposition}[COS commutativity for total-effect indices]
\label{prop:cos_T}
The COS operator $T^{\mathrm{COS}}_B$ commutes with total-effect Sobol index
computation if and only if the pointwise total-effect sensitivity structure
$S_i^T(z)$ is constant over $B$ for all $i$.
\end{proposition}

\begin{proof}
The proof follows the same $L^2$-basis strategy as the proof of
Proposition~\ref{prop:cos}, applied to the total-effect ANOVA components.

Write the complementary ANOVA component at location $z$ as
$f_{\sim i}(\mathbf{x}_{\sim i}, z) := Y(z) - \mathbb{E}[Y(z)\mid X_i] - (\text{interaction terms not involving }X_i)$,
so that $\mathrm{Var}_{X_i}[Y(z)\mid \mathbf{X}_{\sim i}]$ captures the
residual variance of $Y(z)$ not explained by $\mathbf{X}_{\sim i}$.
The total-effect numerator for the block average is
\[
  N_i^T := \mathbb{E}_{\mathbf{X}_{\sim i}}\bigl[\mathrm{Var}_{X_i}[\bar{Y}_B \mid \mathbf{X}_{\sim i}]\bigr]
  = |B|^{-2}\int_B\int_B
    \mathbb{E}_{\mathbf{X}_{\sim i}}\bigl[\mathrm{Cov}_{X_i}[Y(z), Y(z') \mid \mathbf{X}_{\sim i}]\bigr]
    \,d\mu(z)\,d\mu(z').
\]

When $S_i^T(z) = S_i^T$ is constant over $B$ for all $i$, the same
$L^2$-basis argument applies: expand the residual component in an
orthonormal basis of $L^2(\mu_i)$, show that constancy of $S_i^T(z)$
forces the basis coefficients to factor as
$c_k(z) = \alpha_k\sqrt{D_i^T(z)}$ where $D_i^T(z) = S_i^T\,\mathrm{Var}[Y(z)]$,
and deduce that the block-averaged residual is proportional to a single
function of $X_i$ with a spatial scaling factor.  The ratio
$N_i^T / \mathrm{Var}[\bar{Y}_B]$ then equals $S_i^T = \hat{S}_i^{T,B}(Y)$.

When $S_i^T(z)$ varies over $B$, the basis coefficients do not factor,
the block average mixes distinct ``shapes'' of the residual component, and
$S_i^T(\bar{Y}_B) \neq \hat{S}_i^{T,B}(Y)$.  $\square$
\end{proof}

\subsection*{B.3\ \ Proposition~3 for total-effect indices}

\begin{proposition}[Scalar invariance for total-effect indices]
\label{prop:scalar_T}
Let $f: \Omega \to \mathbb{R}$ be a scalar output and $\phi$ strictly monotone.
Then
\[
  \frac{S_i^T(\phi \circ f)}{S_j^T(\phi \circ f)}
  = \frac{S_i^T(f)}{S_j^T(f)}
  \quad \forall\, i \neq j.
\]
Consequently, any two strictly monotone transforms produce identical $D^T$
and $\Delta^T$ on a given scalar benchmark.
\end{proposition}

\begin{proof}
Let $\phi_1, \phi_2$ be strictly monotone.  As in the proof of
Proposition~\ref{prop:scalar}, set $\rho = \phi_1 \circ \phi_2^{-1}$ and
$Z = \phi_2(Y)$, $W = \rho(Z) = \phi_1(Y)$; it suffices to show
$S_i^T(W) = S_i^T(Z)$.

The total-effect index for $W$ is
$S_i^T(W) = \mathbb{E}_{\mathbf{X}_{\sim i}}[\mathrm{Var}_{X_i}[W \mid \mathbf{X}_{\sim i}]] / \mathrm{Var}[W]$.
Since $W = \rho(Z)$ with $\rho$ strictly monotone and $Z$ scalar,
$\sigma(W) = \sigma(Z)$ (by injectivity of $\rho$).  Therefore
$\sigma(W) \vee \sigma(\mathbf{X}_{\sim i}) = \sigma(Z) \vee \sigma(\mathbf{X}_{\sim i})$,
and the conditional distribution of $W$ given $\mathbf{X}_{\sim i} = \mathbf{x}_{\sim i}$
is the image under $\rho$ of the conditional distribution of $Z$ given
$\mathbf{X}_{\sim i} = \mathbf{x}_{\sim i}$.

For each fixed $\mathbf{x}_{\sim i}$, define the conditional distribution
$F_{\mathbf{x}_{\sim i}}(z) := \Pr[Z \leq z \mid \mathbf{X}_{\sim i} = \mathbf{x}_{\sim i}]$.
The conditional variance of $W$ given $\mathbf{X}_{\sim i}$ is
$\mathrm{Var}_{X_i}[\rho(Z) \mid \mathbf{X}_{\sim i} = \mathbf{x}_{\sim i}]$.
Similarly the conditional variance of $Z$ is
$\mathrm{Var}_{X_i}[Z \mid \mathbf{X}_{\sim i} = \mathbf{x}_{\sim i}]$.
Since $\rho$ is monotone but not affine, these conditional variances are
not generally proportional.  However, the key observation is that we are
comparing $W = \phi_1(Y)$ and $Z = \phi_2(Y)$, and both are monotone
functions of the \emph{same} scalar $Y$.  The conditional distribution of $Y$
given $\mathbf{X}_{\sim i} = \mathbf{x}_{\sim i}$ determines both
conditional distributions entirely.

Now apply the copula invariance argument rigorously.  The joint distribution
of $(X_i, Y)$ conditional on $\mathbf{X}_{\sim i}$ has a unique copula
$C_i(\cdot, \cdot \mid \mathbf{x}_{\sim i})$.  The conditional distribution
of $\phi(Y)$ given $(\mathbf{X}_{\sim i}, X_i)$ is the push-forward of the
conditional distribution of $Y$ under $\phi$.  Since $\phi$ is strictly
monotone, the copula of $(X_i, \phi(Y))$ conditional on $\mathbf{X}_{\sim i}$
is identical to the copula of $(X_i, Y)$ conditional on $\mathbf{X}_{\sim i}$
(Sklar's theorem \citep{Sklar1959}; monotone transforms of marginals leave the copula invariant).

The Sobol total-effect index $S_i^T$ depends on the conditional variance of
$Y$ given $\mathbf{X}_{\sim i}$, which is a functional of the conditional
distribution of $Y$ given $\mathbf{X}_{\sim i}$.  When normalised by the
total variance $\mathrm{Var}[Y]$, this ratio is a copula-based quantity:
it measures the fraction of total dependence between $X_i$ and $Y$ that
is not explained by $\mathbf{X}_{\sim i}$.  Since the copula of
$(X_i, \phi_1(Y))$ conditional on $\mathbf{X}_{\sim i}$ equals the copula
of $(X_i, \phi_2(Y))$ conditional on $\mathbf{X}_{\sim i}$, the
normalised conditional variance ratio is identical:
$S_i^T(\phi_1 \circ f) = S_i^T(\phi_2 \circ f)$ for all $i$.

The invariance of $D^T$ and $\Delta^T$ follows immediately, as they are
functions of the index vector $(S_1^T, \ldots, S_d^T)$.  $\square$
\end{proof}

\subsection*{B.4\ \ Numerical verification}

Table~\ref{tab:total_effects} reports $D^T$ and $\Delta^T$ computed from
total-effect indices for representative transforms across all three benchmark
families, using $N = 2048$ (scalar and spatial) or $N = 1024$ (temporal).
First-order values appear in parentheses for direct comparison.

\begin{table}[H]
\centering
\caption{Total-effect index scores $D^T$ and $\Delta^T$ for selected
transforms and benchmarks, with first-order Jansen MC values in parentheses
(all from \texttt{sweep\_results\_v3}, $N=2048$ scalar/spatial, $N=1024$ temporal).
The table includes all three benchmark families (scalar, spatial, temporal)
and both linear and nonlinear transforms.  The Ishigami function shows
the largest gap between $D^T$ and $D$ ($\approx 0.04$) due to its strong
$X_1$--$X_3$ interaction term; near-additive benchmarks (Borehole, Campbell2D,
DampedOscillator) show gaps $\leq 0.04$.  By Proposition~\ref{prop:scalar_T},
all nonlinear transforms on Borehole give identical $(D^T,\Delta^T)$.}
\label{tab:total_effects}
\small
\begin{tabular}{llcc}
\toprule
Transform & Benchmark & $D^T$ ($D$) & $\Delta^T$ ($\Delta$) \\
\midrule
\multicolumn{4}{l}{\textit{Scalar benchmarks}} \\
Arrhenius            & Ishigami     & 0.29 (0.247) & 0.65 (0.612) \\
Log-shift            & Ishigami     & 0.29 (0.247) & 0.65 (0.612) \\
Arrhenius            & Borehole     & 0.62 (0.599) & 0.81 (0.796) \\
Log-shift            & Borehole     & 0.62 (0.599) & 0.81 (0.796) \\
Hill dose-response   & Borehole     & 0.62 (0.599) & 0.81 (0.796) \\
\midrule
\multicolumn{4}{l}{\textit{Spatial benchmarks}} \\
Arrhenius            & Campbell2D   & 0.37 (0.335) & 0.56 (0.527) \\
Exceedance area      & Campbell2D   & 0.29 (0.254) & 0.78 (0.780) \\
Isoline perimeter    & Campbell2D   & 0.32 (0.288) & 0.46 (0.434) \\
\midrule
\multicolumn{4}{l}{\textit{Temporal benchmarks}} \\
Arrhenius            & DampedOsc.   & 0.22 (0.200) & 0.23 (0.196) \\
Bandpass (linear)    & DampedOsc.   & 0.44 (0.431) & 0.73 (0.732) \\
Cumsum (linear)      & DampedOsc.   & 0.07 (0.064) & 0.41 (0.406) \\
\bottomrule
\end{tabular}
\end{table}

The four conclusions from Propositions~\ref{prop:main_T}--\ref{prop:scalar_T}
are all confirmed numerically.  (i)~All nonlinear transforms elevate $D^T$
above zero for every benchmark family, consistent with Proposition~\ref{prop:main_T}.
(ii)~The linear controls (cumsum, bandpass) maintain $D^T \approx D$ for the
same variance-reallocation reasons as the first-order case.  (iii)~Ishigami
shows the largest $D^T > D$ gap ($\approx 0.04$), because the $X_1$--$X_3$ product term
inflates total effects relative to first order, changing the baseline profile
and thereby shifting the threshold-crossing score; Borehole and Campbell2D
show smaller gaps ($\leq 0.04$), consistent with their near-additive structure.
(iv)~Scalar invariance holds for total effects: Ishigami's $D^T$ and
$\Delta^T$ are identical for Arrhenius and log-shift, as Proposition~\ref{prop:scalar_T}
requires.

%% ─────────────────────────────────────────────────────────────────────────────
\section{Benchmark Function Provenance and Applied Literature}
\label{app:benchmarks}
%% ─────────────────────────────────────────────────────────────────────────────

\subsection{Canonical citations and usage references for benchmark functions}

Table~\ref{tab:bench_provenance} lists the canonical source for each benchmark
function used in this study, together with representative sensitivity analysis
papers that have employed it as a test case.

\begin{table}[H]
\centering
\caption{Benchmark function provenance. ``Canonical'' is the primary source
defining the function; ``SA usage'' lists papers that have used it as a
sensitivity analysis test case.}
\label{tab:bench_provenance}
\small
\begin{tabular}{p{2.6cm}p{3.0cm}p{7.5cm}}
\toprule
Benchmark & Canonical source & Representative SA usage \\
\midrule
Ishigami & \citet{Ishigami1990} & \citet{Saltelli2002,Sudret2008,Puy2022} \\
SobolG & \citet{SobolLevitan1999} & \citet{Saltelli2002,Blatman2011,Puy2022} \\
Borehole & \citet{MorrisMitchell1993} & \citet{Saltelli2008,Forrester2008,Puy2022} \\
Piston & \citet{BenAriSteinberg2007} & \citet{Saltelli2008,Puy2022} \\
WingWeight & \citet{Forrester2008} & \citet{Saltelli2008,Puy2022} \\
OTL Circuit & \citet{BenAriSteinberg2007} & \citet{Saltelli2008,Puy2022} \\
Damped Oscillator & \citet{Sudret2008} & \citet{Lamboni2011,Blatman2011} \\
Lotka-Volterra & \citet{Lotka1925,Volterra1926} & \citet{Saltelli2005,Campolongo2011} \\
Campbell2D & \citet{MarrelSpatial2011} & \citet{Marrel2011,Buzzard2012} \\
Campbell3D & \citet{Hettinger2025c3d} & This paper (first use) \\
\bottomrule
\end{tabular}
\end{table}

\subsection{Per-transform applied literature}

Table~\ref{tab:transform_lit} provides representative literature in which each
transform category is applied in practice in combination with sensitivity or
uncertainty analysis.  These references support the ecological validity of the
transform registry and motivate each category's inclusion.

\begin{table}[H]
\centering
\caption{Applied literature for each transform category.}
\label{tab:transform_lit}
\small
\begin{tabular}{p{2.0cm}p{4.5cm}p{7.0cm}}
\toprule
Category & Representative papers & Context \\
\midrule
Environmental & \citet{Vogel1996,Merz2010} & Log-transform of streamflow for
flood frequency; power-law depth-damage in flood loss estimation \\
Environmental & \citet{Thirel2024,BoxCox1964} & Box-Cox and log transforms for
hydrological model calibration; systematic comparison of transform effects \\
Engineering & \citet{Arrhenius1889,Saltelli2005} & Arrhenius rates in atmospheric
chemistry SA; parameter sensitivity of reaction-rate ODEs \\
Engineering & \citet{Weibull1951,SocieMarquis2000,Hill1910} & Weibull failure-probability
in reliability SA; Goodman stress criterion in fatigue; Hill sigmoid in
dose-response models \\
Spatial & \citet{Cressie1993,GotwayYoung2002} & Change-of-support; block-averaging
spatial fields to observational scales \\
Spatial & \citet{Matern1986,MarrelSpatial2011} & Matérn kernel smoothing of
sensitivity fields; spatial Sobol index estimation \\
Temporal & \citet{Lamboni2011,Campbell2006} & Generalised SA indices for dynamic
(ODE) models; functional output sensitivity \\
Temporal & \citet{Sudret2008,Blatman2011} & PCE-based estimation of time-varying
$S_i(t)$; trajectory compression to scalar statistics \\
Statistical & \citet{Iman1979,PlischkeBorgonovo2013} & Rank/PIT transform before
regression-based SA; moment-independent sensitivity from given data \\
Statistical & \citet{Borgonovo2007,Borgonovo2014} & Delta-index as
transformation-invariant alternative; invariance conditions for variance-based indices \\
\bottomrule
\end{tabular}
\end{table}

%% ─────────────────────────────────────────────────────────────────────────────
\clearpage
\bibliographystyle{plainnat}
\bibliography{noncommutativity}
%% ─────────────────────────────────────────────────────────────────────────────

\end{document}
